\section{Identidad con signo y estimación de llegada}\label{app:signed}
Al revelar una reina, la ley uniforme pondera cada hijo por sus completaciones y la ley
adaptativa lo pondera por el predictor. Al comparar ambas a lo largo de una historia hay que
conservar la normalización de cada fila y el valor del predictor al detenerse.
Sea \(\mu\) la ley de la solución uniforme con el orden independiente de filas, y sean \(\nu\)
y \(\nu^\sharp\) las leyes adaptativas de \(L\) y \(L^\sharp\). Dada la fila \(i\), pongamos
\[
 \rho^\sharp(B,i)=\sum_{e\in i}\exp(L^\sharp(B+e)-L^\sharp(B)),\qquad
 A^\sharp=\sum_{t<\tau}\log\rho^\sharp(B_t,i_t).
\]
Aquí \(t\) cuenta los pasos de revelado en la suma y todos los hijos disponibles, incluidos los de salida,
intervienen en la normalización. Sea \(g^\sharp=\log Q-L^\sharp\). En la llegada conservemos
sólo las restricciones de pares de un cuadrado físico \(C\) sobre la misma máscara de torres.
Denotemos por \(\gamma_C\) el logaritmo de la probabilidad de evitación bajo torres menos el
cambio correspondiente en el predictor. Sean \(\Gamma=\sum_C\gamma_C\) y
\(g^\flat=g^\sharp-\Gamma\). Los factores del orden de filas son iguales en las dos historias.
Las demás probabilidades se telescopan, pues la ley uniforme asigna a cada prefijo una masa
proporcional a su número de completaciones. Al invertir en \(F_b\), donde dicho número es
positivo y la historia tiene masa positiva bajo ambas leyes, se obtiene la identidad exacta
\begin{equation}\label{eq:signed-history}
 e^{-g^\sharp(B_n)}\nu^\sharp(F_b)
 =\mathbb E_\mu[\exp\{-A^\sharp-\Gamma-g^\flat\};F_b].
\end{equation}
En adelante \(\mathbb E[X;E]=\mathbb E[X\mathbf1_E]\). Para estimar la identidad necesitamos
la consistencia por fila y la difusión en los tableros que entran en la normalización.
Precisamos esta última propiedad para el hijo físico, también cuando es el primero que sale.
\begin{lema}[Escalado después de una jugada]\label{lem:hijo-fisico}
Sea un padre regular de tamaño \(m\) cuyo punto físico \(p^{Q}\) es positivo en su soporte
binario, tiene márgenes de fila y columna uno, \(p^{Q}_{ic}\le C/m\) y brecha de incidencia
al menos \(\gamma>0\). Fijada una casilla legal \(e\), fórmese el menor de torres \(M_e\)
y prohíbase cualquier subconjunto de las demás casillas de las dos diagonales de \(e\).
Para \(m\) suficientemente grande, todas esas máscaras de tamaño \(m-1\) tienen escalado en
una misma clase de difusión y brecha, con constantes que sólo dependen de \(C,\gamma\).
En particular esto vale para el hijo completo \(M_e^{\rm fis}\) en la primera salida.
\end{lema}
\begin{proof}
Eliminamos una columna en cada fila retenida y hasta una casilla por cada diagonal.
La pérdida en cada fila es a lo sumo \(3C/m\), y lo mismo vale por columna. Sea \(R\) la
restricción de \(p^Q\) después de las eliminaciones y sea \(\mathcal W\) el lado con menos vértices
de un corte del grafo bipartito restante. En el padre, Rayleigh sobre el indicador centrado
da capacidad al menos \(\gamma |\mathcal W|/2\), pues \(|\mathcal W|\le m\).
Queda en \(R\) una capacidad de al menos \((\gamma/2-3C/m)|\mathcal W|\).
Oriente las aristas de filas a columnas. Los déficits de fila y los déficits de columna
con signo opuesto tienen suma cero y cada uno tiene módulo a lo sumo \(3C/m\).
Con \(\lambda=(3C/m)/(\gamma/2-3C/m)<1\), cada corte soporta la demanda neta
con capacidades simétricas \(\lambda R\). El criterio de flujo factible~\cite{gale} da una corrección
firmada
y una matriz bistocástica \((1-\lambda)R\le p'\le(1+\lambda)R\).
Su conductancia es al menos \((1-\lambda)(\gamma/2-3C/m)\). La desigualdad de Cheeger~\cite{lsok} da una
brecha
positiva uniforme para el laplaciano normalizado.

Falta transferir estas cotas de una tabla factible al escalado del soporte binario.
En tamaño \(k\), sea \(a\) bistocástica, positiva en el soporte, con \(a_{ic}\le C'/k\)
y brecha \(\gamma'\). El máximo de entropía en ese soporte es interior por la factibilidad
estricta. Sus ecuaciones dan el escalado \(D\) con \(\log(kD_{ic})=u_i-u_c\).
Por los márgenes, \(\sum a_{ic}\log(kD_{ic})=\sum D_{ic}\log(kD_{ic})\ge0\).
La parte positiva de la primera suma es a lo sumo \(\sum a_{ic}kD_{ic}\le C'k\).
Al restar una mediana de los potenciales e integrar los cortes de nivel se obtiene
\[
 \frac{\gamma'}2\sum_z|u_z|\le\sum_{ic}a_{ic}|u_i-u_c|\le2C'k .
\]
Fije \(T>8(C')^2/\gamma'\). Cada línea tiene al menos \(k/C'\) vecinos, y en todo el grafo
menos de \(k/(2C')\) vértices tienen \(|u_z|>T\). La ecuación de fila
\(1=e^{u_i}\sum_{c\sim i} e^{-u_c}/k\) da \(e^{u_i}\le2C'e^T\), y la de columna da
\(e^{-u_c}\le2C'e^T\). Las mismas ecuaciones con las cotas superiores dan cotas inferiores
recíprocas. Así \(d_-/k\le D_{ic}\le d_+/k\) en el soporte y
\(D\ge(d_-/C')a\) conserva una brecha positiva por Rayleigh.
Se aplica esto a \(p'\), con \(k=m-1\). Para la máscara del padre se aplica directamente a
\(p^Q\), también en la llegada larga regular. Los márgenes del predictor~\cite{pc} se obtienen
hasta la parada por el Corolario~4.12 y la clase interior de~\cite[def.~2.2 y §§4--5]{tc}.
Si la brecha se enuncia con pesos de respuesta, su comparabilidad con \(p^Q\) en la rama
de vacancias positivas da la hipótesis de Rayleigh.
\end{proof}

El coeficiente \(\coefT=\mathfrak F-1/12\) de (T) depende de las masas marcadas de fila y
columna \(s,t\) y del peso del cruce \(z\). Las expresiones integradas del
apéndice~\ref{app:rook} son holomorfas y acotadas en discos con radios de clase.
En \(s=0\) o \(t=0\) los términos visibles son independientes de \(z\).
Al derivar en \(z\), ambas caras se anulan. Si \(M_0\) acota \(\coefT\) en un
polidisco de radio \(r\) de clase en torno al punto de borrado en \(z\), Cauchy en
\(z\) acota \(|\partial_z\coefT|\) por \(2M_0/r\) en el disco reducido. Schwarz en las
masas \(s,t\) introduce los factores \(|s|/r,|t|/r\). Por tanto
\(\partial_z\coefT=st\,O(1)\) uniformemente, también al integrar en \(z\) a lo largo
del segmento de borrado. Las dos masas valen \(1/k\); en tamaño \(k\) esto da
\(O(k^{-3})\) para el cambio de \(\coefT/k\) por casilla vetada. Retirar una fila y una columna
cambia las masas en \(O(k^{-1})\) y el tamaño inverso en \(O(k^{-2})\), con coste
\(O(k^{-2})\) en \(\coefT/k\). Hay a lo sumo \(2m\) casillas vetadas. El lema proporciona
los mismos márgenes en todos los pasos. Sumando los cambios y usando el resto \(O(m^{-2})\)
de (T) en los dos extremos, obtenemos para todos los hijos del padre regular
\begin{equation}\label{eq:incremento-hijo-fisico}
 |g_T(M_e^{\rm fis})-g_T(M)|\le C/m^2,\qquad |g_T(M_e^{\rm fis})|\le C/(m-1).
\end{equation}


\paragraph{El signo de la salida.}
La cota \eqref{eq:incremento-hijo-fisico} y la consistencia del predictor ordinario dan
\(\lvert A^\sharp\rvert\le C_b/n\) hasta la parada, al sumar \(O(m^{-2})\) en \(m>j\).
En \(F_b\) la difusión da \(|\Gamma|\le C/j\). Pongamos \(U_0=A^\sharp+\Gamma\).
El Lema~\ref{lem:llegada-integrada} acota el coste de reemplazar \(e^{-g^\flat}\) por uno en
\eqref{eq:signed-history}, pues \(e^{-U_0}\le e^{C_b/n}\).
El resto de Taylor de \(e^{-U_0}\) es \(O_b(n^{-2})\).
La cota previa \(R_Q(n)=O(n^{-1})\) del Lema~\ref{lem:previa}
y (T) dan \(g^\sharp(B_n)=O(n^{-1})\), con el mismo orden de resto para su exponencial.
Por tanto, con \(\limsup_n|\epsilon_n(b)|\le C_\delta b\),
\[
 n g^\sharp(B_n)\nu^\sharp(F_b)
 =n\mathbb E_\mu[U_0;F_b]+n\{\mu(F_b^c)-\nu^\sharp(F_b^c)\}+\epsilon_n(b).
\]
El signo se obtiene de \(\nu^\sharp(F_b)-\mu(F_b)=\mu(F_b^c)-\nu^\sharp(F_b^c)\).

El cociente de probabilidades de hijo bajo \(\nu^\sharp\) y \(\nu\), dada la fila, es
\(e^{g_T(M_e^{\rm fis})}/\mathbb E_{\nu,i}e^{g_T(M_{e'}^{\rm fis})}\).
Todos los hijos del padre regular cumplen \(|g_T|\le C/(m-1)\), incluido el de primera
salida. El logaritmo del cociente está entre \(-2C/(m-1)\) y \(2C/(m-1)\), pues el
denominador promedia sobre todos los hijos. El producto detenido queda acotado por
\(\exp(2C\sum_{m=j+1}^n(m-1)^{-1})=\exp(O_{b,\delta}(1))\). Con \(X=Qe^{-L}\),
la densidad de \(\mu\) respecto de \(\nu\) es \(Xe^{A-R_Q(n)}\), donde \(A\) suma los
logaritmos de las normalizaciones de fila ordinarias. En el tramo macroscópico la probabilidad
de salida bajo \(\nu\) es exponencialmente pequeña por~\cite[cor.~4.12]{tc}, y el coste
ponderado por \(e^AX\) se acota por
\(C\exp(-c_\delta n+C_\delta n^{2/3}+O_\delta(n^{-1}))\).
Para cada banda diádica de la parada con suelo \(j\), de escala \(m\ge j/2\), el coste con
peso \(1+X\) es
\(C e^{-(c_*-2\beta)m}\), con \(c_*>2\beta\).
El Lema~\ref{lem:prescripcion-larga} acota la llegada larga en \(O(e^{-2j})\) con ese peso.
Al sobrevivir el tramo macroscópico se cumple \(|A|\le C_\delta/\lceil\delta n\rceil+C/j\).
El Lema~\ref{lem:positividad-llegada} da \(F_b^c\subset H_j^c\). La suma por bandas y la
cota macroscópica controlan \(\nu(H_j^c)\) y \(\mu(F_b^c)\) a orden exponencial.
El producto detenido da ese orden para \(\nu^\sharp(F_b^c)\). Retirar el indicador en la media
de \(A^\sharp\) cuesta entonces \(o_b(1)\) tras multiplicar por \(n\). Se sigue que
\begin{equation}\label{eq:signed-mean}
 \limsup_n\left|n g^\sharp(B_n)-n\mathbb E_\mu A^\sharp
                 -n\mathbb E_\mu[\Gamma;F_b]\right|\le C_\delta b .
\end{equation}
\section{Suma del resto en la llegada}\label{app:arrival}
Un solo cuadrado puede contribuir a la escala \(1/n\) que interesa en \eqref{eq:signed-mean}.
Al sustraer \(\Gamma\), cada soporte que sobrevive debe satisfacer las restricciones geométricas
de dos cuadrados. Ese coste adicional permite sumar los errores con una cota que decrece con \(b\).
Veámoslo primero en la llegada histórica corta \(H_j\), bajo \(\nu\). Para un conjunto
de casillas \(Y\), activamos sólo las restricciones diagonales soportadas en \(Y\), manteniendo
la máscara de torres completa. Llamemos \(e(Y)\) al logaritmo de la probabilidad de evitación
menos el cambio de predictor, y \(e_C(Y)\) a la misma expresión usando sólo los pares de \(C\).
El coeficiente de un soporte \(U\) es \(\sum_{Y\subseteq U}(-1)^{|U|-|Y|}e(Y)\).
Una casilla aislada se cancela al emparejar subconjuntos que difieren sólo en él. Si todos los
pares pertenecen a un cuadrado, el coeficiente se cancela con el de \(e_C\).
Si intervienen dos cuadrados distintos, ningún \(e_C\) contribuye: las casillas tendrían que
estar en los cuatro vértices de \(C\), cuyas únicas parejas diagonales son las dos de \(C\).
El resto suma por tanto los coeficientes originales que carecen de casillas aisladas y contienen
al menos dos cuadrados distintos. Lo sustraído tiene grado a lo sumo cuatro.

Sean \(v\) el número de casillas, \(h\) el total de etiquetas distintas de filas y columnas,
\(D\) el número de coordenadas libres y \(p\) el exponente del soporte. Prescribir las
\(h\) líneas da el factor \((j/n)^h\), junto con \(\delta^{-2v}\exp(C_Lv^2/j)\) en la
cota de prescripción. Para cada tipo, el número de posiciones es de orden \(n^D\).
Escribamos \(d_*=h-D\), \(q=p-D-1\). Dos cuadrados distintos imponen dos ecuaciones
lineales no proporcionales sobre las coordenadas de fila y columna. Así \(d_*\ge2\);
la cota del soporte da \(q\ge0\). El factor escalado es
\[
 n\,n^D j^{-p}(j/n)^h=(j/n)^{d_*-1}j^{-q}
 \le (b/\delta_*)\delta_*^{d_*-1}j^{-q}
 \qquad(0<b\le\delta_*).
\]
Extraemos \(b/\delta_*\) antes de sumar. Para \(v\le4\) el número de tipos tiene una cota
absoluta. Para los tipos no colineales el recuento ponderado usa \(d_*-1+q\ge(v-4)/3\).
De ahí que para \(5\le v\le N=\lceil4\log n/\log\rho\rceil\), fijado \(\rho\ge2\), la cota
tras sumar los tipos de ese tamaño sea
\[
 (b/\delta_*) C_L(C_1\delta^{-2})^v v^{d_{\rm pol}}\delta_*^{(v-4)/3}.
\]
Se elige \(\delta_*\) de manera que \(C_1\delta^{-2}\delta_*^{1/3}<1\).
En un tipo colineal \(v\ge3\) se tiene \(p=v+2,h=2v\), con a lo sumo \(4n\,n^v/v!\)
soportes. Resulta \(C(C_2\delta^{-2})^v v^2(j/n)^{v-2}\); se reduce \(\delta_*\)
una vez para sumar también esta cota. La sustitución de átomos hasta grado \(N\) cuesta
\(O(j^{-8})\) en valor, y la norma global acota la cola por \(O_\rho(n^{-2})\).
Ambos errores desaparecen tras escalar para cada \(b\) fijo. Así
\(\limsup_n n\mathbb E_\nu[|g^\flat|;H_j]\le C_\delta b\), con \(C_\delta\)
independiente de \(b\).

La cota obtenida controla una media de \(|g^\flat|\) bajo \(\nu\). En \eqref{eq:signed-history}
interviene además la exponencial bajo \(\mu\). El factor \(X\) de la densidad permite estimar
ambas con la misma cota, pues en la llegada se tiene
\[
 X\,|e^{-g^\flat}-1|=e^{g_T+\Gamma}|1-e^{g^\flat}|.
\]
En \(H_j\) la parte positiva del error diagonal tiene una cota de clase,
\(|g_T|+|\Gamma|\le C/j\) y \(|A-R_Q(n)|=O_{b,\delta}(1/n)\).
Se aplica \(|e^y-1|\le e^{y^+}|y|\) a la densidad \(Xe^{A-R_Q(n)}\).
En \(F_b\setminus H_j\) la probabilidad ponderada es \(O(e^{-2j})\).
Para \(0<X\le1\) vale \(X|\log X|\le1/e\); por encima de uno se usa
\(X|\log X|\le X(C_0+\beta j)\) en la clase exterior. El aporte de este suceso al lado
izquierdo de \eqref{eq:arrival-estimate} antes de escalar es \(O((1+j)e^{-2j})\).
Para obtener \eqref{eq:arrival-estimate} fijamos primero \(\delta\), los márgenes, \(\rho\) y
las constantes de clase, después \(\delta_*\) y por último el corte. El umbral de \(n\) puede
depender de \(b\). La constante \(C_\delta\) sirve para todos esos cortes, como exige la
comparación final de dos cortes en el apéndice~\ref{app:closing}.

\section{La media interior a corte positivo}\label{app:interior}
Para hallar la contribución interior a \eqref{eq:signed-mean} fijamos el corte \(b>0\).
Cuando quedan \(m>\lfloor bn\rfloor=j\) filas, escribimos \(x=m/n\).
La elección de una casilla modifica el predictor en orden \(1/m\) después de descontar
\(\log p_c\). Para identificar lo que queda en orden \(1/m^2\) empezamos en un padre regular,
antes de promediar sobre las historias.

\subsection{La inserción de una casilla}
Sean \(c\) la casilla elegida, \(\mathcal J=\familias(c)\) sus cuatro líneas, \(I\) las líneas retenidas,
\(F\) las casillas disponibles que tocan \(\mathcal J\), incluido \(c\), y \(E'\) las que sobreviven.
Para cada casilla \(i\), sea \(a_i\) el vector de incidencia de sus cuatro líneas.
Con \(h_f=1-\sum_{i\in f}p_i\) en las diagonales, pongamos
\begin{align*}
 S_0&=\sum_i p_i a_i a_i^{\mathsf T}+\operatorname{diag}(0_R,0_C,h_D,h_A),
 & R&=S_0^+,\\
 t_{ic}&=a_i^{\mathsf T}Ra_c,
 & w_i&=6p_i+f_i,\qquad f_i=\sum_{\{u,v\}\subset \familias(i)}R_{uv}.
\end{align*}
Las sumas sobre pares no ordenados usan dos líneas distintas. En \(E'\), escribamos
\(y_i=p_i^{(c)}-p_i\), donde \(p^{(c)}\) es el punto estacionario del hijo, y
\(q_i=y_i/p_i+t_{ic}\). La respuesta del hijo da \(q_i=O_b(m^{-2})\).
El potencial completo del prefactor es
\(V=\ell+\tfrac12\log\det S_{0,\rm gf}+3\sum_i p_i^2\). En los soportes conexos de
la clase, \(L=H+\ell+\tfrac12\log(2\pi)\) y \(\ell\) incluye el cofactor del predictor.
En \(\det S_{0,\rm gf}\) se usa el mismo calibre en padre e hijo.
La corrección lineal y su centrado se leen en
\[
 \delta_c^0=-w_c+\sum_i p_i t_{ic}w_i,\qquad
 \sum_{c\in r}p_ca_c=S_0 e_r,\qquad a_i^{\mathsf T}RS_0 e_r=(a_i)_r .
\]
La última igualdad usa que \(a_i\) es ortogonal al modo nulo fila--columna.
Por tanto \(\sum_{c\in r}p_c\delta_c^0=0\) para cada fila \(r\).

La consistencia por inserción~\cite[Teorema 7.1 y Corolario 7.2]{pc} da
\(L(B+c)-L(B)=\log[p_c(1+\delta_c)]+r_c^{\rm PC}\), con
\(\sum_{c\in r}p_c\delta_c=0\), \(\delta_c=O_b(m^{-1})\) y \(r_c^{\rm PC}=O_b(m^{-2})\).
Al desarrollar la exponencial en esos pesos de masa uno y después el logaritmo se obtiene
\(\log\rho(B,r)=\sum_{c\in r}p_c r_c^{\rm PC}+O_b(m^{-3})\).
El cambio de pesos en el término lineal cuesta \(O_b(m^{-3})\).

\paragraph{Qué términos deja el cambio de soporte.}
Descompongamos \(S_0\) en bloques \(B_0,C_0,C_0^{\mathsf T},D_0\) según \(I,\mathcal J\).
Sea \(U_0=\sum_{i\in F}p_i a_{i,I}a_{i,I}^{\mathsf T}\),
\(\mathcal M=B_0-U_0\), \(d_0=C_0\mathbf1\), \(\beta=Ra_c\),
\(z=\beta_I\), \(t=\beta_{\mathcal J}-\mathbf1\) y \(r=d_0+\mathcal M z=-C_0t-U_0z\).
El mínimo de la energía cuadrática del hijo es \(Q_*=d_0^{\mathsf T}\mathcal M^+d_0\).
Al insertar \(d_0=r-\mathcal Mz\) y usar \(S_0\beta=a_c\) se obtiene
\[
 Q_*=a_c^{\mathsf T}Ra_c-4+12p_c+E_Q,\qquad
 E_Q=z^{\mathsf T}U_0z-t^{\mathsf T}D_0t+r^{\mathsf T}\mathcal M^+r .
\]
En efecto \(D_0=(1-p_c)I+p_c\mathbf1\mathbf1^{\mathsf T}\) y
\(\mathbf1^{\mathsf T}D_0\mathbf1=4+12p_c\). Se tiene \(r=O_b(m^{-2})\) por entrada
y \(r^{\mathsf T}\mathcal M^+r=O_b(m^{-3})\) por la brecha y el número \(O_b(m)\) de líneas.

Sea \(\gamma_f=h_f^{(c)}-h_f\) el cambio de vacancia en una diagonal retenida.
Taylor de la entropía del apéndice~\ref{app:predictor}, con los pesos del hijo, da
\begin{align*}
 \Delta H-\log p_c&=\tfrac32\sum_{F\setminus\{c\}}p_i^2-\tfrac12 Q_*+A_3+O_b(m^{-3}),\\
 A_3&=\sum_{F\setminus\{c\}}p_i^3+\tfrac16\sum_{E'}p_i(y_i/p_i)^3
 +\tfrac32\sum_{E'}\frac{y_i^2}{1-p_i}+\tfrac16\sum_{f\in I_D}\gamma_f^3/h_f^2 .
\end{align*}
Aquí \(I_D\) recorre ambas familias diagonales retenidas. Los términos cuárticos cuestan
\(O_b(m^{-3})\) por \(\sum p_i=m\), \(p_i=O_b(m^{-1})\) y
\(|y_i/p_i|+|\gamma_f/h_f|=O_b(m^{-1})\). La ortogonalidad en el mínimo deja un error
cuadrático en las correcciones relativas \(O_b(m^{-2})\) a la respuesta lineal.

Para el cofactor escribamos \(\Pi_m=vv^{\mathsf T}/(2m)\), donde \(v\) vale \(1,-1,0,0\)
en las cuatro familias. Sean \(H_*=(R+\Pi_m)_{II}\), \(C_{\mathcal J}=D_0+(\Pi_m)_{\mathcal J\mathcal J}\),
\(V_{\mathcal J}=C_0+(\Pi_m)_{I\mathcal J}\) y \(\Pi'=v_Iv_I^{\mathsf T}/[2(m-1)]\).
En \(I\), \(S_*\) tiene entrada \(-p_i\) en cada par ordenado de líneas distintas
de \(i\in F\setminus\{c\}\). En forma análoga \(Y\) tiene entrada \(y_i\) en los pares de
\(i\in E'\). Sus restantes entradas son cero. El complemento de Schur da
\begin{align*}
 D_*&=S_*+Y+\Pi'-(\Pi_m)_{II}+V_{\mathcal J} C_{\mathcal J}^{-1}V_{\mathcal J}^{\mathsf T},\\
 Q_c&=\tfrac{m^2}4\{\operatorname{tr}[(H_*D_*)^2]-\operatorname{tr}(S_*^2)\},\\
 E_R&=\operatorname{tr}[(C_{\mathcal J}-I)(H_{\mathcal J}-I)],\qquad H_{\mathcal J}=(R+\Pi_m)_{\mathcal
 J\mathcal J}.
\end{align*}
Para la normalización del calibre sean \(S'_0\) la matriz del hijo y
\(\Delta_0=S'_0-(S_0)_{II}\). Su contribución se obtiene de
\[
 E_g=\log[m/(m-1)]-\log\det C_{\mathcal J}
   +\operatorname{tr}[(\Pi_m)_{II}\Delta_0]
   +\operatorname{tr}[H_*(\Pi'-(\Pi_m)_{II})].
\]
Como \(v_{\mathcal J}\perp\mathbf1\), \(\det C_{\mathcal J}=(1+3p_c)(1-p_c)^2(1-p_c+1/m)\).
Los logaritmos dan \(m^{-2}-p_c/m+6p_c^2+O_b(m^{-3})\). Las dos trazas valen
\(-(1-p_c)/m\) y \((m-1)/m^2+v_{\mathcal J}^{\mathsf T}R_{\mathcal J\mathcal J}v_{\mathcal J}/[2m(m-1)]\),
 respectivamente,
por \(S_0v=Rv=0\). Como \(R_{\mathcal J\mathcal J}=I+O_b(m^{-1})\), queda \(E_g=m^{-2}+6p_c^2+O_b(m^{-3})\).
El cubo del logaritmo del determinante deja \(C^{(3)}=-\operatorname{tr}(S_*^3)/6\).
En efecto, las filas de \(S_*\) tienen suma absoluta \(O_b(m^{-1})\);
las entradas de \(D_*-S_*\) son \(O_b(m^{-2})\) y las de \(H_*-I\) son \(O_b(m^{-1})\).
En la traza cúbica, un factor de \(H_*D_*-S_*\) cuesta así \(O_b(m^{-3})\) en total.
Las potencias de orden cuatro o superior tienen la misma cota por la norma de filas.
La traza pura \(\operatorname{tr}(S_*^2)/4\) aporta
\(3\sum_{F\setminus\{c\}}p_i^2/2\), pues hay seis pares ordenados en cada casilla retirada.

En la parte lineal de Schur, sea \(f_i^I\) la restricción de \(f_i\) a pares retenidos.
Las cuatro entradas diagonales de \(S_0R=I-\Pi_m\) sobre \(\mathcal J\) dan
\[
 \sum_{i\in F}p_i(f_i-f_i^I)
 =-\operatorname{tr}[(R-I)_{\mathcal J\mathcal J}]-1/m-p_cf_c .
\]
Al juntar la traza lineal, \(-Q_*/2\), las dos copias de
\(3\sum_{F\setminus\{c\}}p_i^2/2\) y \(-3\Delta\sum p_i^2\), los términos con \(w_i\)
son \(\sum_F p_iw_i-w_c-\sum_{E'}y_iw_i\) y queda \(p_cf_c-3p_c^2\).
Restar \(\delta_c^0\) convierte las dos sumas en las de la última línea de la fórmula
\begin{align*}
 L(B+c)-L(B)-\log p_c-\delta_c^0
   &=\Delta_c V+P_c+O_b(m^{-3}),\\
 P_c&=-E_g/2+Q_c/m^2+C^{(3)}+A_3-3\sum_{E'}y_i^2\\
 &\quad-3p_c^2+p_cf_c-(E_Q+E_R)/2\\
 &\quad+\sum_F p_i(1-t_{ic})w_i-\sum_{E'}p_iq_iw_i .
\end{align*}


Puesto que \(\delta_c-\delta_c^0=O_b(m^{-2})\), desarrollar \(\log(1+\delta_c)\) da
\[
 r_c^{\rm PC}=\Delta_c V+P_c+\tfrac12(\delta_c^0)^2
                   +(\delta_c^0-\delta_c)+O_b(m^{-3}).
\]
El último paréntesis desaparece exactamente al promediar en la fila. Los errores cúbicos
acumulados cuestan \(O_b(n\sum_{m>j}m^{-3})=O_b(n^{-1})\). La tabla identifica qué
parte de \(m^2[P_c+(\delta_c^0)^2/2]\) origina cada contribución al perfil.
\begin{center}\small
\begin{tabular}{ll}
\toprule
Término de inserción a escala \(m^2\) & Contribución al perfil\\ \midrule
\(-m^2 E_g/2\)& \(G_{\rm cof}\)\\
\(m^2 A_3\)&\(f2E+f2z+f2F\)\\
\(m^2 C^{(3)}\)&\(f2F+C_{3,\rm cuad}\)\\
\(-m^2(E_Q+E_R)/2\)&\(f3F+f3c\)\\
\(m^2(-3p_c^2+p_c f_c)\)&\(f4=f4_{\rm expl}+f4_{L0}\)\\
\(m^2\sum_F p_i(1-t_{ic})w_i\)&\(f5c+f5F\)\\
\(-m^2\sum_{E'}p_iq_iw_i\)&\(f6\)\\
\(m^2(\delta_c^0)^2/2\)&\(f7\)\\
primera parte de \(Q_c\)&\(Q_1+\Delta_1\)\\
segunda parte de \(Q_c\), con \(-3m^2\sum_{E'}y_i^2\)&\(Q_2+\Delta_2\)\\
\bottomrule
\end{tabular}
\end{center}
\subsection{Las dieciséis contracciones del perfil}\label{sec:contracciones}
Denotamos por \(\langle g\rangle_i=\int_{[0,1]^2}\nu_1(i)g(i)\,di\) la media de área y
por \(\langle\cdot\rangle_c\) la misma media en el centro. En la fibra \(\lambda\)
de \(c\), \(d\omega_\lambda(i)=\eta_\lambda\nu_1(i)\,dy\), con \(y\) la coordenada libre
de fila o columna y, en las diagonales, la coordenada de fila. Para una función \(g\),
\(\int_F g=\sum_{\lambda\in \familias(c)}\int_{i\in\lambda}g(i,\lambda)\,d\omega_\lambda(i)\).
La medida \(d\rho\) de una línea libre es \(dy\) en
cada familia axial y \(U_f\,dy/(2x)\) en cada hoja de una familia diagonal.
El factor \(1/2\) ya realiza la media en la hoja de esa línea libre. Los demás índices
enteros libres aportan sus propias clases de paridad, según su parametrización. En una fibra
axial el desplazamiento tiene dos paridades de igual peso. En una diagonal el desplazamiento
mantiene el color. Un cruce diagonal--antidiagonal sólo existe en hojas compatibles.

\paragraph{Qué disponibilidad lleva una intersección.}
En una configuración fijamos las líneas de sus casillas marcadas y las líneas que se integran
con \(d\rho\). Sea \(\mathcal A\) ese conjunto. A cada otra línea \(\ell\) de una configuración
finita se le asigna una variable de Bernoulli independiente \(X_\ell\) de media \(U_\ell\).
En \(\mathcal A\) se pone \(X_\ell=1\). Sean \(\familias(z)\) las cuatro líneas de \(z\),
\(Y_z=\prod_{\ell\in \familias(z)}X_\ell\) y \(\pesodirecto(z)=W(z)Y_z\). Un cruce ausente da
 \(\pesodirecto=0\).
Así el peso de varias casillas directas \(Z\) se obtiene de
\[
 \psi(Z)=\mathbb E\prod_{z\in Z}Y_z
 =\prod_{\ell\in(\bigcup_{z\in Z}\familias(z))\setminus\mathcal A}U_\ell .
\]
El mismo indicador se usa en todo el monomio. Para el núcleo \(L_0\) y la corrección
\(\theta\) de la sección~\ref{sec:profile-data}, la respuesta de dos líneas es
\[
 E(a,b)=
 \begin{cases}
 (L_0)_{aa}+\theta_a,&a=b,\\
 (L_0)_{ab}-\pesodirecto(a\cap b),&a\ne b .
 \end{cases}
\]
Una casilla marcada aporta \(\pesodirecto(i)=W(i)\). Abreviamos
\(z_f=\sum_{j\in \familias(c)}E(f,j)\), \(T_i=\sum_{a\in \familias(i)}z_a\),
\(t_j=z_j\) para \(j\in \familias(c)\), y \(w(i)=\sum_{\{a,b\}\subset \familias(i)}(L_0)_{ab}\).
En particular \(t_j\) y \(w(i)\) en una casilla marcada son deterministas.
Para \(i\) que sólo comparte \(\lambda\) con \(c\), sea
\(\sigma_i=\sum_{a\in \familias(i)\setminus\{\lambda\}}z_a\).

\paragraph{Cubo, energía y corrección de fila.}
Las contribuciones del cubo de la acción son
\begin{align*}
 f2E&=\left\langle-\tfrac16\mathbb E T_i^3+\tfrac32 W(i)\mathbb E T_i^2\right\rangle_{i,c},\\
 f2z&=-\tfrac16\left\langle\int_{D\cup A}(u_f/U_f)\,\mathbb E z_f^3\,d\rho(f)\right\rangle_c,\\
 f2F&=\left\langle\int_F W(i)^2\right\rangle_c .
\end{align*}
En \(f2E\) se marcan \(i,c\), independientes en área. En \(f2z\) se marcan \(c,f\).
Para verificar los momentos en estos dos casos, sumemos sobre los pares de líneas que definen
\(T_i\) o \(z_f\). Sean \(A_1,A_2,A_3\) las sumas respectivas de \(L_0-K\),
\(K(W_{ab}-K)\) y \(-K(W_{ab}-K)(W_{ab}-2K)\), con \(W_{ab}=W(a\cap b)\) o cero.
En posición general las transversales de cruces distintos son distintas. El segundo momento es
\(A_1^2+A_2\) y el tercero \(A_1^3+3A_1A_2+A_3\), por los momentos centrados de cada Bernoulli.
El factor \(u_f/U_f\) es la vacancia del punto estacionario en la línea retenida.

Con \(I_2=\langle W(c)^2\rangle_c\) y \(\Xi=\langle W(c)w(c)\rangle_c\), se obtiene
\[
 G_{\rm cof}=-\tfrac12-3I_2,\qquad f4=-9I_2+\Xi,\qquad f5c=-3\Xi .
\]
En efecto \(mf_c\) tiende al factor \(w(c)-6W(c)\), y \(1-t_{cc}=-3+O_b(m^{-1})\).
Las otras contracciones del centro y de la fibra son
\begin{align*}
 f3F&=-\tfrac12\left\langle\int_F\mathbb E\sigma_i^2\right\rangle_c,\\
 f3c&=\tfrac12\left\langle\sum_{j\in \familias(c)}t_j^2
 -\sum_{j,h\in \familias(c)}[W(c)(1-\delta_{jh})+\Pi_{jh}]
                         [E(h,j)+\Pi_{hj}]\right\rangle_c,\\
 f5F&=-\left\langle\int_F w(i)\mathbb E T_i\right\rangle_c,\\
 f7&=\tfrac12\left\langle\left[-w(c)+\int_F w(i)
                  +\left\langle w(i)\mathbb E T_i\right\rangle_i\right]^2\right\rangle_c .
\end{align*}
Aquí \(\Pi_{jh}=v_jv_h/2\) es el núcleo del perfil. El primer sumando de \(f3c\)
procede de \(t^{\mathsf T}D_0t\) y el segundo de \(E_R\).
En \(f3F\), \(i,c\) están marcadas en la misma fibra. Sean \(Z\) los cruces existentes
\(a\cap j\), con \(a\in \familias(i)\setminus\{\lambda\}\) y \(j\in \familias(c)\setminus\{\lambda\}\),
de familias distintas. En posición general son casillas distintas. Escribiendo
\(\psi_z=\psi(\{z\})\), el signo negativo de cada parte directa desaparece en la covarianza,
\begin{align*}
 \mathbb E\sigma_i^2={}&(\mathbb E\sigma_i)^2+\sum_{z\in Z}W(z)^2\psi_z(1-\psi_z)\\
 &+\sum_{\substack{z,z'\in Z\\z\ne z'}}W(z)W(z')
       [\psi(\{z,z'\})-\psi_z\psi_{z'}].
\end{align*}
Por ejemplo, en coordenadas relativas a \(c=(0,0)\), si \(i=(0,s)\), los cruces
\((s,s)\) y \((s,0)\) comparten la fila \(s\). Cuando existen dentro del tablero, su peso
conjunto es \(\psi_z\psi_{z'}/U_R\). En \(f7\) las dos copias de la suma tienen índices
libres distintos. Tras fijar \(c\), sus coincidencias accidentales tienen medida nula.
Así se cuadran las integrales de la respuesta centrada antes de promediar en \(c\).

\paragraph{Triángulos de líneas en la traza cúbica.}
Tres líneas del producto \(S_*^3\) concurren en una casilla borrada o cortan tres casillas
distintas que completan un cuadrado con \(c\). En el primer caso sus seis órdenes cancelan
el divisor \(6\), y el resultado es otra copia de \(f2F\).
Para el segundo, tomemos la base \(((1,0),(0,1))\) o \(((1,1),(1,-1))\) y multipliquemos
cada vector por un signo independiente, obteniendo \(v,w\). Sean
\(Z=\{c,c+hv,c+hw,c+h(v+w)\}\), \(h>0\), y \(n_{RC}\) el número de líneas axiales de \(Z\).
Las tres coordenadas libres son las de \(c\) y el paso primitivo \(h\). El resultado es
\[
 C_{3,\rm cuad}=\sum_{v,w}\int_{\{(c,h):\,Z\subset[0,1]^2\}}
  x^{n_{RC}-3}\prod_{z\in Z}W(z)
  \prod_{\ell\in(\bigcup_{z\in Z}\familias(z))\cap(D\cup A)}U_\ell\,dc\,dh .
\]
Los cuatro factores \(W\) incluyen el peso del centro. En la cuadratura el lado cero
lleva la mitad del valor límite del integrando. En la integral tiene medida cero.

\paragraph{Respuesta de segundo orden en el exterior.}
La fila \(f6\) requiere \(m^2q_i\), pues la respuesta de primer orden en \(y_i/p_i\)
se cancela con \(t_{ic}\). Para verificar la fuente de segundo orden, pongamos
\(X=(a_{i,I}^{\mathsf T})_{i\in E'}\), \(D_p=\operatorname{diag}(p_i)_{i\in E'}\).
Las ecuaciones estacionarias restadas dan \(y_i/p_i=(X\eta)_i+(B_1)_i\), donde
\(\eta_f=\gamma_f/h_f\) en las diagonales y las otras componentes son multiplicadores
de fila y columna. Con \(\theta_i=y_i/p_i\), la diferencia exacta de los logaritmos
de las ecuaciones de estacionariedad da
\[
 (B_1)_i=\theta_i-\log(1+\theta_i)-3\log(1-y_i/(1-p_i))
 +\sum_{f\in \familias(i)\cap(D\cup A)}[\log(1+\eta_f)-\eta_f].
\]
Las ecuaciones de masa dan
\(\mathcal M\eta=d_0-X^{\mathsf T}D_pB_1\). Por tanto
\[
 q=B_1+X\mathcal M^+(r-X^{\mathsf T}D_p B_1).
\]
Taylor en las casillas y las vacancias da \(m^2B_1=L+O_b(m^{-1})\), con
\[
 L_i=\tfrac12(mt_{ic})^2-\frac{3mp_i(mt_{ic})}{1-p_i}
       -\tfrac12\sum_{f\in \familias(i)\cap(D\cup A)}(mz_f^{\rm fis})^2,\qquad
 z^{\rm fis}=(Ra_c)_I .
\]
Aquí se han usado \(y_i/p_i=-t_{ic}+O_b(m^{-2})\) y
\(\gamma_f/h_f=-z_f^{\rm fis}+O_b(m^{-2})\).
En el complemento del modo nulo, la identidad \(S_0R=I-\Pi_m\) por bloques da
\(\|\mathcal M^+-R_{II}\|_\infty=O_b(m^{-1})\) en norma de suma por fila,
pues las líneas retenidas pierden masa \(O_b(m^{-1})\) y la inversa está acotada.
Así \(m^2q=L+XR_{II}(m^2r-X^{\mathsf T}D_pL)+O_b(m^{-1})\).

Expresamos la contracción en los mismos núcleos del perfil. Fijemos el centro \(c\) y sea
\(\Lambda(i)=T_i^2/2-3W(i)T_i-\tfrac12\sum_{f\in \familias(i)\cap(D\cup A)}z_f^2\).
Para una línea retenida \(u\), escribamos
\begin{align*}
 \omega(u)={}&-\sum_{j\in \familias(c)} \pesodirecto(u\cap j)
     \left[t_j+\sum_{a\in \familias(u\cap j)\setminus\{j\}}z_a\right]
      -\int_{i'\in u}\mathbb E\Lambda(i')\,d\omega_u(i'),\\
 q(i)={}&\Lambda(i)+\sum_{u\in \familias(i)}
 \left[\omega(u)+\int E(u,u')\omega(u')\,d\rho(u')\right],\\
 f6={}&-\left\langle w(i)\mathbb E q(i)\right\rangle_{i,c}.
\end{align*}
Se omiten los cruces inexistentes en la primera suma. El término de fibra integrado en
\(\omega\) es la media con \(i'\) marcada y \(c\) fijo. En los productos de la suma sobre
cruces se usa \(\psi\) para todo el monomio. Las dos partes de \(\omega\) provienen de
\(m^2r=-m^2(C_0t+U_0z)\) y de \(-X^{\mathsf T}D_pL\). La identidad y la parte integral
de \(R_{II}\) producen los dos términos entre corchetes de \(q(i)\). La media de área
lleva el signo de \(-m^2\sum_{E'}p_iq_iw_i\). El Lema~\ref{lem:f6} demuestra su
reordenación por trazas, incluido el recorrido de la figura~\ref{fig:rectangulos}.

\begin{lema}[Reordenación de la contribución \(f6\)]\label{lem:f6}
En la malla de lado \(s\), sean \(\mathcal I{}\) la incidencia casilla--línea de la
sección~\ref{sec:rectangulos}, \(\beta_e=\nu_1(e)/s^2\) y
\(N_v=\mathcal I{}^{\mathsf T}\operatorname{diag}(v)\mathcal I{}\). Los productos de vectores de casillas
son puntuales. Con \(K,L_0,\rho,w\) de la sección~\ref{sec:profile-data}, definamos
\begin{align*}
 D&=L_0-K,& V_{ij}&=K_{ij}(W_{ij}-K_{ij}),&
 \Psi&=\rho^{-1}\mathcal I{}^{\mathsf T}(\beta w),\\
 Z&=\Psi+D\rho\Psi,& h&=w-\mathcal I{}Z,& \alpha&=\beta h,\\
 a&=\mathcal I{}^{\mathsf T}\alpha,& b_\beta&=\mathcal I{}^{\mathsf T}\beta,&
 M&=\mathcal I{}D\mathcal I{}^{\mathsf T}.
\end{align*}
Aquí \(W_{ij}\) es el peso del cruce o cero si éste no existe.
Escribamos \(U(e,j)=\sum_{u\in \familias(e)} K_{uj}-3W(e)\) y
\(t(c,e;j)=M_{ce}+\theta_j+U(c,j)+U(e,j)\).
La cuadratura del perfil físico, con el valor genérico en coincidencias accidentales,
da exactamente \(f6_s=-E_1+E_2\), donde
\begin{align*}
 E_1={}&\tfrac12\operatorname{tr}(N_\beta D N_\alpha D)
       +\tfrac12 b_\beta^{\mathsf T}Va-3b_\beta^{\mathsf T}D\mathcal I{}^{\mathsf T}(\alpha W)\\
     &-\tfrac12\sum_{f\in D_{\rm obl}}a_f
         [(D N_\beta D)_{ff}+(Vb_\beta)_f],\\
 E_2={}&\sum_c\beta_c\sum_{j\in \familias(c)}\sum_{e\in j}
     \frac{\eta_j\nu_1(e)}s\,t(c,e;j)
     \sum_{u\in \familias(e)\setminus\{j\}}Z_u .
\end{align*}
En esta fórmula \(D_{\rm obl}\) reúne las dos familias oblicuas.
\end{lema}
\begin{proof}
Fijado \(c\), la respuesta de segundo orden consta de la fuente local \(\Lambda\) y
la fuente de líneas \(\omega\). Al sumar la respuesta contra \(\beta_e w(e)\), cada
\(u\in \familias(e)\) recibe \([\mathcal I{}^{\mathsf T}(\beta w)]_u=(\rho\Psi)_u\).
La parte integral de la respuesta añade \(\rho D\rho\Psi\) por simetría. La masa
total es, pues, \(\rho Z\). La fuente negativa de \(\omega\) integra \(\Lambda\) en
la fibra de \(u\). En la malla \(\rho_u\eta_u=1/s\), por lo que esa suma sustrae
\(\beta_e(\mathcal I{}Z)_e\) del peso de \(\Lambda(e;c)\). Queda la masa \(\alpha_e\).

Para dos marcas genéricas \(c,e\), la media de \(T_e\) es \(M_{ce}\).
Cada cruce directo tiene varianza \(V_{uj}\), pues su indicador al cuadrado es el propio
indicador. Los cruces distintos tienen transversales distintas en posición genérica.
Así \(\mathbb E T_e^2=M_{ce}^2+\sum_{u\in \familias(e),j\in \familias(c)}V_{uj}\).
Para \(f\in \familias(e)\) la misma cuenta da
\(\mathbb E z_f^2=(\sum_{j\in \familias(c)}D_{fj})^2+\sum_{j\in \familias(c)}V_{fj}\).
Al sumar \(\beta_c\alpha_e\mathbb E\Lambda(e;c)\), el cuadrado \(M_{ce}^2\)
da la traza de \(N_\beta D N_\alpha D\) por multiplicación de índices.
La varianza da \(b_\beta^{\mathsf T}Va\), y el momento lineal con \(W(e)\)
da \(b_\beta^{\mathsf T}D\mathcal I{}^{\mathsf T}(\alpha W)\).
Los cuadrados de \(z_f\) dan \((DN_\beta D)_{ff}+(Vb_\beta)_f\), uno por línea
oblicua de \(e\). Los coeficientes de \(\Lambda\) producen \(E_1\).
En \(N_v\) se incluyen las entradas diagonales \(\sum_{e\in u}v_e\), también cuando
el recorrido vuelve por la misma línea.

En la suma de cruces de \(\omega\), fijemos \(e=u\cap j\), con \(j\in \familias(c)\).
El peso directo por la medida de \(u\) es \(\rho_uK_{uj}=\eta_j\nu_1(e)/s\).
Condicionar el cruce a estar disponible marca entonces \(c,e\) en la fibra \(j\).
En la respuesta media \(M_{ce}\) se repone la diagonal \(\theta_j\).
Los tres cruces con \(j\) en cada extremo llevan ahora el peso cierto \(W\) de la
casilla marcada. Cambiar cada antiguo término \(-K\) a \(-W\) añade \(U(c,j)+U(e,j)\).
Queda \(t(c,e;j)\), multiplicado por \(Z_u\) para cada transversal \(u\ne j\) de \(e\).
La fuente de cruces tiene signo negativo en \(\omega\), y la definición de \(f6\)
tiene otro signo negativo. De ahí \(+E_2\) y \(-E_1\). Esta igualdad de sumas
finitas deja el error espacial de la extrapolación en el estado \status{N}.
\end{proof}

\paragraph{Las dos trazas del término mixto.}
Pongamos \(D=H_*D_*-S_*\). La ciclicidad de la traza da la identidad en el padre
\[
 Q_c=\tfrac12\sum_{u,v}(S_*)_{uv}(m^2D)_{vu}
       +\tfrac1{4m^2}\sum_{u,v}(m^2D)_{uv}(m^2D)_{vu}.
\]
Cada casilla de \(F\setminus\{c\}\) tiene tres líneas retenidas, que dan seis pares
ordenados en la primera suma. En la segunda se integran dos líneas libres.
Para expresar \(m^2D\) en el perfil escribamos \(\mathcal K(v,w)=E(v,w)+\Pi_{vw}\),
\(\mathcal V(v,j)=\pesodirecto(v\cap j)+\Pi_{vj}\), \(j\in \familias(c)\). Para líneas retenidas \(v,u\),
sea \(\mathcal Y(v,u)=-\pesodirecto(z)T_z\) si se cortan en \(z\in E'\), y cero en otro caso.
En esta fórmula \(T_z\) usa los mismos indicadores que \(\pesodirecto(z)\).
Finalmente, para \(w\ne u\) retenidas, sea \(\mathcal S(w,u)=-\pesodirecto(i')\) cuando
\(w\cap u=i'\in F\setminus\{c\}\), y cero en los demás pares. Definamos
\begin{align*}
 \mathcal B(v,u)&=\mathcal Y(v,u)+\Pi_{vu}
                    +\sum_{j\in \familias(c)}\mathcal V(v,j)\mathcal V(u,j),\\
 \mathcal E_2(v,u)&=\mathcal B(v,u)+\sum_w\mathcal K(v,w)\mathcal S(w,u)
                     +\int\mathcal K(v,w)\mathcal B(w,u)\,d\rho(w).
\end{align*}
Estas tres piezas proceden, con \(B_*=D_*-S_*\) y \(K_*=H_*-I\), de
\(m^2D=m^2B_*+(mK_*)(mS_*)+m^{-1}(mK_*)(m^2B_*)\).
En la suma discreta se toma cada \(i'=u\cap j\), \(j\in \familias(c)\), y cada una de
sus dos líneas transversales \(w\). Su disponibilidad está en \(\pesodirecto(i')\).
En la integral sí se marca \(w\), pues su disponibilidad ya está en \(d\rho(w)\).

Escribamos \(R(i)=\{(u,v):u,v\in \familias(i)\setminus\{\lambda\},\ u\ne v\}\) para una
casilla de la fibra \(\lambda\). Aplicando \(\psi\) a los monomios se obtienen
\begin{align*}
 L_1&=-\tfrac12\left\langle\int_F
          \sum_{(u,v)\in R(i)}\mathbb E\mathcal E_2(v,u)\right\rangle_c,\\
 L_2^{\rm int}&=\tfrac14\left\langle\iint
 \mathbb E\big[\mathcal E_2(u,v)\mathcal E_2(v,u)-\mathcal Y(u,v)\mathcal Y(v,u)\big]
                 \,d\rho(u)\,d\rho(v)\right\rangle_c .
\end{align*}
Las dos copias de una línea integrada internamente tienen variables libres separadas.
El bloque \(\mathcal Y\mathcal Y\) se sustrae antes de integrar. En el tablero cada
\(i\in E'\) aporta doce pares ordenados al bloque \(YY\), de modo que su contribución
es \(+3m^2\sum_{E'}y_i^2\), exactamente la opuesta a la de \(P_c\).
También se puede comprobar en el continuo. Para cada par de familias, la medida de los
cruces válidos por su disponibilidad es \(U_DU_A\,di\), con las hojas y el jacobiano del
cruce incluidos. Multiplicando por \(W(i)^2\) resulta \(\nu_1(i)W(i)\).

\paragraph{Cómo extraer \(\Delta_1\) y \(\Delta_2\).}
Las filas tabuladas son \(Q_1=L_1-\Delta_1\) y \(Q_2=L_2^{\rm int}-\Delta_2\).
Damos las correcciones por separado para poder reconstruir cada entrada. En todas las
sumas siguientes un cruce debe existir en el cuadrado y en la clase de paridad indicada.
Una igualdad de líneas se exige como identidad en los parámetros libres de esa clase.
Las coincidencias accidentales forman conjuntos de medida nula.

Para \(\Delta_1\), fijemos \(c,i,\lambda,(u,v)\in R(i)\) como en \(L_1\), y una
línea libre \(w\) integrada con \(d\rho(w)\). Pongamos \(z=v\cap w,\ p=w\cap u\).
Para cada \(a\in \familias(p)\setminus\{w,u\}\) y \(j\in \familias(c)\), pongamos \(t=a\cap j\).
Se marcan \(\familias(c)\cup \familias(i)\cup\{w\}\).
El monomio directo de \(\mathcal K(v,w)\mathcal Y(w,u)\) aporta la diferencia
\[
 \Delta_1=\tfrac12\left\langle\int_F\sum_{(u,v)\in R(i)}
 \int\sum_{a,j} W(z)W(p)W(t)
 \big[\psi(\{z,p,t\})-\psi(\{z\})\psi(\{p,t\})\big]\,d\rho(w)\right\rangle_c .
\]
Se toman los tres cruces distintos, con \(p\in E'\). La diferencia cuenta la transversal
común de \(z,t\) sólo una vez. Si es \(\ell\), el corchete vale
\(\psi(\{z\})\psi(\{p,t\})(U_\ell^{-1}-1)\).
Las tres partes directas llevan signo negativo en el producto de la primera traza
antes de aplicar su prefactor \(-1/2\).
Por ejemplo, en coordenadas relativas, para \(c=(0,0)\), \(i=(0,s)\), \(u=C_s\),
\(v=D_{-s}\), \(w=R_h\), \(a=A_{h+s}\), \(j=R_0\), resultan \(z=(h,h+s)\),
\(p=(h,s)\), \(t=(0,h+s)\). La columna \(h+s\) aparece en \(z\) y \(t\).

Para \(\Delta_2\), marquemos \(\familias(c)\cup\{u,v\}\), con \(u,v\) libres. Cuando existe el cruce,
el peso de sus transversales es
\[
 K_{ab}=W(a\cap b)\prod_{\ell\in \familias(a\cap b)\setminus\{a,b\}}U_\ell .
\]
Para cada \(j\in \familias(c)\) tomamos \(p_2=v\cap j\),
\(w_2\in \familias(p_2)\setminus\{v,j\}\), \(t_2=u\cap w_2\), \(k_2=K_{vj}K_{uw_2}\).
Los dos patrones posibles para los otros dos vértices dan
\[
\begin{array}{c|c|c|c}
 &p_1,\ t_1 &k_1&\text{multiplicidad }s\\ \hline
 \mathcal K\mathcal S\ \text{por}\ \mathcal K\mathcal S&
 u\cap j,\ v\cap w_1 & K_{uj}K_{vw_1}&1\\
 \mathcal Y\ \text{por}\ \mathcal K\mathcal S&
 v\cap u,\ a\cap j & K_{vu}K_{aj}&2
\end{array}
\]
En la primera fila \(w_1\) recorre \(\familias(p_1)\setminus\{u,j\}\). En la segunda \(a\)
recorre \(\familias(p_1)\setminus\{v,u\}\). En cada patrón se conservan cuatro vértices distintos.
Para cada pareja de cruces que define \(k_d\), \(d=1,2\), reunamos las transversales de
sus dos cruces, es decir, las líneas cuyos factores \(U\) aparecen en \(k_d\).
Si los dos conjuntos comparten una línea \(\ell\) fuera de las marcadas, se incluye el
sumando siguiente. Los cuatro cruces forman entonces un cuadrado y la transversal es única.
\[
 \Delta_2=\tfrac14\left\langle\iint\sum_{\text{patrones admitidos}}
          s\,k_1k_2(U_\ell^{-1}-1)\,d\rho(u)\,d\rho(v)\right\rangle_c .
\]
Se puede obtener la lista desarrollando los productos de \(\mathcal E_2\).
Un peso \(\pesodirecto\) aporta las cuatro coordenadas de su casilla menos las líneas ya marcadas.
Para comparar dos etiquetas de una misma familia se restan sus expresiones afines en
las coordenadas libres. Sólo una diferencia idénticamente nula aporta el factor
\(U_\ell^{-1}-1\). Para hacer la comprobación basta comparar el término constante y
los coeficientes de cada parámetro en las cuatro etiquetas \(r,q,r-q,r+q\). En
\(\Delta_1\) se obtienen 24 elecciones de familias y 80 al distinguir las clases admisibles.
En \(\Delta_2\) quedan 24 elecciones en cada fila de la tabla. Por ejemplo, en la primera,
con \(j=R_0,u=A_s,v=D_t,w_1=C_s,w_2=C_{-t}\), los dos vértices exteriores
\((s+t,s),(s+t,-t)\) comparten la fila \(s+t\). En la parte mixta, intercambiar las copias da el mismo
integrando con \(u,v\) permutadas, de donde procede el \(2\) de la tabla.
Finalmente,
\begin{align*}
 d(x)={}&G_{\rm cof}+f2E+f2z+2f2F+f3F+f3c+f4+f5c+f5F\\
        &+f6+f7+C_{3,\rm cuad}+Q_1+\Delta_1+Q_2+\Delta_2 .
\end{align*}


\subsection{Límite uniforme de las medias y suma sobre los niveles}
La regularidad antes de la parada da \(p_i\le C_b/m\), una vacancia diagonal
minorada y una brecha de filas y columnas~\cite[\S7]{pc}.
En consecuencia \(c_b I\le S_0+\Pi_m\le 4I\) para algún \(c_b>0\).
Para ver la cota inferior, la forma cuadrática es la suma de
\(\sum_i p_i(a_i^{\mathsf T}y)^2\) y \(\sum_f h_f y_f^2\), más la proyección de calibre.
El segundo término controla las componentes diagonales y la brecha controla las axiales
ortogonales al calibre. La cota superior sigue de la diagonal unidad de \(S_0\), de que su
masa fuera de la diagonal por fila es a lo sumo tres y de que la proyección actúa en su núcleo.

Sea \(A_m=S_0+\Pi_m-I\), \(M=(I+A_m)^{-1}=R+\Pi_m\) y \(r_b=1-c_b/4<1\).
El polinomio \(P_N=\tfrac14\sum_{h=0}^N[I-(I+A_m)/4]^h\) satisface
\(\|M-P_N\|_{\rm op}\le r_b^{N+1}/c_b\). Para controlar una entrada multiplicada
por \(m\) usamos \(M=I-A_m+A_mMA_m\). La suma absoluta de una fila de \(A_m\) es
a lo sumo cuatro y su máximo es \(O_b(m^{-1})\), luego
\(\|A_me_u\|_2=O_b(m^{-1/2})\) para todo \(u\). Así
\[
 \sup_{u,v}m\,|[A_m(M-P_N)A_m]_{uv}|\le C_b r_b^{N+1}.
\]
Esto acota el error en cada entrada de \(m(R-I)\) al sustituir \(M\) por \(P_N\)
en la identidad anterior. Las entradas originales y sus truncaciones están acotadas.

A \(N\) fijo, las potencias de \(S_0-I\) son sumas sobre sucesiones de líneas adyacentes.
Las inserciones de \(\Pi_m\) se simplifican usando \((S_0-I)\Pi_m=-\Pi_m\).
Un producto de recorridos es una suma finita de tipos de incidencia. Si dos etiquetas afines
son idénticas en un tipo, se cuenta la línea una vez. Una coincidencia adicional impone
una ecuación no constante y deja \(O_{b,N}(n^{d-1})\) posiciones entre las \(O(n^d)\)
de un tipo con \(d\) parámetros. En la diagonal, la ida y vuelta de longitud dos fuerza
la misma casilla en ambos pasos. Su contribución es \(3m\sum_{i\in u}p_i^2\), pues hay tres
elecciones de la otra línea. Su exceso sobre el producto de las medias da precisamente
\(\theta_u=3\eta_u s_u-\int K_{uv}^2\,d\rho(v)\). Los cruces entre diagonales se
cuentan en las clases de paridad de la sección~\ref{sec:contracciones}.

A cada tipo fijo se aplica el seguimiento de líneas hasta el corte auxiliar \(b/4\)
de~\cite[\S4]{tc}. A tolerancia fija \(\eta>0\), fuera de una excepción de probabilidad
exponencial, \(mp_i=W(i)(1+O_b(\eta+m^{-1}))\) en los sitios vivos y
\(h_f=u_f/U_f+O_b(\eta+m^{-1})\). Para un número fijo de marcas el riesgo en un paso es
la suma de sus riesgos individuales salvo \(O_{b,N}(n^{-2})\) por dobles ocupaciones.
Al multiplicar estas supervivencias hasta el nivel \(m\), una marca axial aporta \(x\)
y una diagonal aporta \(U_f(x)\), pues \(\partial_x\log U_f=(1-u_f)/U_f\).
La aproximación es uniforme en las marcas y en \(j\le m\le n\). Se toma primero \(n\to\infty\)
y luego \(\eta\downarrow0\). Las sumas normalizadas de cada clase son sumas de Riemann
sobre politopos fijos. Los factores del perfil son continuos en \(x\) con módulo común en
\([b,1]\) en cada pieza. Por tanto los truncados convergen uniformemente en los niveles
a integrales continuas de los recorridos.

Para aplicar el mismo argumento a \(mV\), sea \(\kappa_i=a_i^{\mathsf T}Ra_i-4\).
La expansión del prefactor~\cite[\S7]{pc} da, con la última suma sobre todas las líneas,
\[
 V=-14\sum_i p_i^3-\tfrac32\sum_i p_i^2\kappa_i
       +\tfrac94\sum_u\left(\sum_{i\in u}p_i^2\right)^2+O_b(m^{-2}).
\]
Las doce entradas fuera de la diagonal en \(m\kappa_i\) aportan \(-12W(i)\) en la
casilla marcada. Combinadas con \(-14\) dan \(4\langle W^2\rangle_i\).
Las partes largas y la diagonal corregida dan los sumandos con \(L_0,\theta\) de
la sección~\ref{sec:profile-data}. El cuadrado de la suma de fibra da su último término.

En las contracciones, el centro tiene masa uno, el exterior masa a lo sumo uno y las
cuatro fibras masa a lo sumo cuatro. La medida de las líneas libres tiene masa \(\le6/b\).
Los pasos de cuadrados tienen masa total \(\le8/b\). Al desarrollar la diferencia de productos,
la cota de cada entrada se multiplica sólo por masas y factores acotados para \(b\) fijo.
Así la cola del grado \(N\) es \(\le C_b r_b^{N+1}\), uniforme en el nivel tanto para la media
física como para su integral límite. El resto de la inserción en la escala \(m^2\) es
\(O_b(m^{-1})\). Tomar primero \(n\to\infty\) y después \(N\to\infty\) da
\begin{align*}
 \sup_{j<m\le n}\Big|
 \mathbb E_\nu\!\left[m^2\sum_{c\in i_m}p_c
       \big(P_c+(\delta_c^0)^2/2\big);H_j\right]-d(m/n)\Big|&\longrightarrow0,\\
 \sup_{j\le m\le n}\left|\mathbb E_\nu[mV(B_m);H_j]-\mathcal V(m/n)\right|
  &\longrightarrow0 .
\end{align*}
Aquí \(i_m\) es la fila revelada y \(H_j\) la llegada regular corta. Para promediar la fila
se usa primero el indicador de supervivencia anterior a su elección. Dado el padre, la fila
es uniforme y la media se convierte en \(m^{-1}\sum_c p_c(\cdot)\).
Cambiar ese indicador por \(\mathbf1_{H_j}\) en los observables acotados cuesta
\(O_b(\nu(H_j^c))\). La segunda fórmula usa todas las familias en el potencial.
Este argumento identifica una media sobre padres y da continuidad del límite en \([b,1]\).
La evaluación decimal de esas integrales por mallas requiere además una cota de su error espacial.

\paragraph{Los extremos y el cambio de ley.}
El error uniforme \(\epsilon_n(b)\) de la primera media contribuye a lo sumo
\(\epsilon_n(b)n\sum_{m>j}m^{-2}=o_b(1)\). El término de \(P_c+(\delta_c^0)^2/2\)
da por tanto \(\int_b^1d(x)x^{-2}\,dx\). Para \(\Delta_cV\), la probabilidad de un hijo
dada la fila bajo \(\nu\) es \(p_c(1+O_b(m^{-1}))\), y \(|\Delta_cV|\le C_bm^{-2}\).
La sustitución en la media cuesta \(O_b(m^{-3})\). Los incrementos centrados hasta la parada
forman una martingala de segundo momento \(O_b(\sum_{m>j}m^{-4})\).
Si \(M_\tau\) es su suma detenida, las diferencias son ortogonales en \(L^2(\nu)\).
Así \(n|\mathbb E_\nu[M_\tau;H_j]|\le n(\mathbb E_\nu M_\tau^2)^{1/2}\nu(H_j)^{1/2}\).
Como \(\sum_{m>j}m^{-4}=O(j^{-3})\), el coste a escala \(n\) es \(O_b(n^{-1/2})\).
En cada historia de \(H_j\) los incrementos realizados suman \(V(B_j)-V(B_n)\).
La segunda media uniforme da \(\mathcal V(b)/b-\mathcal V(1)\).

Hasta la parada, \(n|A|\le C_b\) para el predictor ordinario. Por el Lema~\ref{lem:previa},
\(n|R_Q(n)|\) también está acotado. En \(H_j\) ponemos
\(s=A-R_Q(n)+g_T+\Gamma\). Allí \(|s|=O_b(n^{-1})\) y la densidad inversa de historias
es \(e^{-s-g^\flat}\). Por \eqref{eq:arrival-estimate},
\[
 \int_{H_j}|d\mu-d\nu|
 \le\mathbb E_\mu[|1-e^{-s}|+e^{-s}|1-e^{-g^\flat}|;H_j]=O_b(n^{-1}) .
\]
Con las estimaciones de salida del apéndice~\ref{app:signed} se sigue que
\(n\mathbb E_\mu A=n\mathbb E_\nu[A;H_j]+o_b(1)\).
Para las torres pongamos \(u_e=g_T(M_e^{\rm fis})-g_T(M)\).
La máscara \(M_e^{\rm fis}\) se obtiene del menor simple \(M_e\) aplicando los dos vetos
diagonales del hijo. El cambio es \(O_b(m^{-2})\), incluido el paso que entra en la salida.
Así \(\log\rho^\sharp-\log\rho=\log\mathbb E_{\nu,i}e^{u_e}
 =\mathbb E_{\nu,i}u_e+O_b(m^{-4})\).
La densidad detenida y las cotas ponderadas de salida dan
\(D(\mu^\tau\Vert\nu^\tau)=\mathbb E_\mu[A+\log X]-R_Q(n)=O_b(1)\).
Sean \(d_m\) las entropías condicionales de las dos leyes de hijo dada la fila,
iguales a cero después de la parada. Las filas tienen la misma ley uniforme.
La descomposición de la entropía relativa en sus incrementos condicionales da
\(\sum_m\mathbb E_\mu d_m=D(\mu^\tau\Vert\nu^\tau)\). La desigualdad de
Pinsker~\cite{ct} acota la diferencia de las dos medias de \(u_e\) por \(2C_bm^{-2}\sqrt{d_m/2}\).
Por Jensen y Cauchy--Schwarz, la suma de estas diferencias a escala \(n\) cuesta
\(O_b(n(\sum_{m>j}m^{-4})^{1/2}D(\mu^\tau\Vert\nu^\tau)^{1/2})=O_b(n^{-1/2})\).
Los incrementos realizados de torres se telescopan y añaden
\((F(b)-1/12)/b-(F(1)-1/12)\) a los extremos anteriores. En conclusión, a cada corte fijo,
\[
 \lim_{n\to\infty} n\mathbb E_\mu A^\sharp
 =\int_b^1\frac{d(x)}{x^2}\,dx+\frac{G(b)}b-G(1)
 =\int_b^1\frac{\numerr(x)}{x^2}\,dx .
\]


\section{Las tres medias de frontera}\label{app:boundary}
En \(\Gamma\) se ha aislado cada cuadrado, pero sus dos pares diagonales pueden estar disponibles
a la vez. Para guardar el signo y contar una sola vez su interacción, empezamos por la
corrección de un par y después reunimos los dos pares en el cuadrado. Aquí \(p_e\) designa
la ocupación estacionaria de \emph{torres} en la máscara binaria.
Para un par diagonal disponible \(P=\{e,f\}\), sea \(\pi^{\rm ex}_P\) su probabilidad conjunta en
una colocación uniforme de torres. Sea \(\pi^T_P=\exp[L_T(M_P)-L_T(M)]\), donde \(M_P\) elimina sus
cuatro extremos de fila y columna. Con \(\mathsf G(e,f)\) la respuesta cruzada del punto estacionario
de torres y \(\chi_e\) su corrección relativa de una casilla, definimos
\[
 \pi_{1,P}=p_ep_f(1-\mathsf G(e,f)+\chi_e+\chi_f),\qquad
 X_P=p_ep_f(p_e+p_f)\mathsf G(e,f).
\]
Para un par solo la corrección del cuadrado es
\(\log(1-\pi^{\rm ex}_P)-\log(1-\pi_{1,P})-X_P+O(j^{-5})\).
Las probabilidades son \(O(j^{-2})\) y su diferencia es \(O(j^{-4})\); al integrar la derivada
logarítmica el error de linealización es \(O(j^{-6})\). El cuadrado completo da las correcciones
de sus dos pares y un término adicional \(-3\prod_{s\in\square}p_s+O(j^{-5})\).
Hay a lo sumo \(j^2(j-1)\) pares, pues cada casilla tiene a lo sumo \(2(j-1)\) vecinas y cada
par tiene dos extremos. Así, la siguiente expresión \(B_n(b)\) difiere de \(n\mathbb E_\mu[\Gamma;F_b]\)
en \(O(n/j^2)=o_b(1)\):
\begin{align*}
 B_n(b)={}&-n\mathbb E_\mu[\textstyle\sum_P(\pi^{\rm ex}_P-\pi^T_P);F_b]\\
           &-n\mathbb E_\mu[\textstyle\sum_P(\pi^T_P-\pi_{1,P}+X_P);F_b]
            -3n\mathbb E_\mu[\textstyle\sum_\square\prod_{s\in\square}p_s;F_b].
\end{align*}
La inserción de \(\pi^T_P\) en esta descomposición tiene dos funciones. La primera media mide
el error del predictor de torres en el par. La segunda conserva el error de inserción
para un par marcado. Para promediarlo hay que saber cuáles de los otros dos vértices siguen
disponibles. Sean \(u_f\) las vacancias finales de las dos familias diagonales y
\(U_f(b)=u_f+b(1-u_f)\). Sean \(v_b=U_D(b)U_A(b)\), \(p_b\) su escalado
bistocástico~\cite{sk} y
\(W_b^0=p_b/v_b\). Un cuadrado de coordenadas mínimas \((r,c)\) y lado \(h\) se integra sobre
\(\Omega=\{0<h<1,\ 0<r,c<1-h\}\).
Con \(d_s=U_D(b,r-c+sh)\), \(a_s=U_A(b,r+c+sh)\), la disponibilidad de su par orientado
tiene factor \(B_o\). Los otros dos vértices tienen indicadores \(I_l I_u,I_l I_v\).
Sus cuatro probabilidades conjuntas y el factor del par marcado son los siguientes.
\[
\begin{array}{c|cccc}
o&B_o&l&u&v\\\hline
 D&d_0a_0a_2&a_1&d_{-1}&d_{+1}\\
 A&a_1d_{-1}d_{+1}&d_0&a_0&a_2
\end{array}
\qquad
\begin{array}{c|c}
s&q_{o,s}\\\hline
00&1-l+l(1-u)(1-v)\\
10&lu(1-v)\\
01&l(1-u)v\\
11&luv
\end{array}
\]
La submedida de cada estado es \(b^4 B_o q_{o,s}\,dr\,dc\,dh\). Las cuatro etiquetas vivas son
dos filas y dos columnas. Una vacancia diagonal repetida conserva su mismo indicador.
Por orientación hay \(\sum_{h=1}^{n-1}(n-h)^2\) posiciones.
La normalización de Riemann aporta \(n^3\), con cuatro factores \(j^{-1}\) y un \(n\) exterior.
Por tanto,
\(n^4 b^4/j^4\to1\). Al insertar los dos pesos marcados resulta
\(w_{o,b}=B_oW_b^0(e)W_b^0(f)=p_b(e)p_b(f)/U_o(b)\).
Se conservan incluso los estados de masa nula sin dividir por \(q_{o,s}\).

Para la primera media, al integrar \(p_b\) sobre cada línea obtenemos su masa \(m_f\). Sean
\[
 W_b=\sum_{f=D,A}\int\frac{m_f^2}{2U_f},\qquad
 \gamma_b=\frac{m_D}{U_D}+\frac{m_A}{U_A},\qquad
 \theta_b(r)=\int p_b(r,c)\gamma_b(r,c)\,dc .
\]
Eliminar los cuatro extremos en la expansión de torres da el término de tamaño \(W_b/6\), el
de normalización \(-2F(b)W_b\) y la variación de las medidas de fila y columna \(-dF_2(b)\).
Aquí \(dF_2(b)\) es la derivada del funcional de torres a núcleo fijo al multiplicar ambas
medidas uniformes por \(1+\varepsilon(2W_b-\theta_b)\), con la simetría de columna.
La dirección tiene masa cero por \(\int\theta_b=2W_b\). Resulta
\(P_T=\lim_{b\downarrow0}[(1/6-2F(b))W_b-dF_2(b)]\).
Para la tercera media, el cuadrado usa seis etiquetas diagonales distintas. Dos aparecen dos
veces entre los ocho factores de vacancia de \(\prod_s W_b^0(s)\). Al multiplicar por las seis
vacancias distintas y tomar \(p_0=\nu_1\) queda exactamente
\[
 -S=-3\int_\Omega
       \frac{\prod_{s\in\square}\nu_1(s)}
            {u_D(r-c)\,u_A(r+c+h)}\,dr\,dc\,dh .
\]
Así quedan identificadas la primera y la tercera medias. La segunda conserva dos datos de la
máscara en el cuadrado marcado, los indicadores de los cruces. El apéndice~\ref{app:weights}
calcula el error del par con esos dos indicadores fijos.

\section{El funcional de una línea en el cuadrado marcado}\label{app:weights}
La segunda media de frontera compara la inserción conjunta con su aproximación al primer
orden. Insertar \(e\) cambia el peso de \(f\) porque elimina su propia fila y columna.
Ordenamos \(e,f\) por fila creciente y conservamos los dos indicadores cruzados \(X,Y\).
Aquí \(M_e\) es el menor simple de torres, sin la fila ni la columna de \(e\). Las primas
se evalúan en ese menor. Sean \(t=p'_f/p_f-1\), \(G_{ef}=\mathsf G(e,f)\) y
\(D=\hat c_e+\hat c_f-G_{ef}\). El producto exacto de los cocientes de inserción da
\[
 e^{u}=\frac{\pi^T_{ef}}{p_ep_f}
      =(1+\hat c_e)(1+\hat c'_f)(1+t)\exp(r_e+r'_f).
\]
La estimación de incrementos de torres \((A^+)\) del Apéndice~\ref{app:rook} es
\(r_i=-\mathsf E_i+O(j^{-3})\), también en el menor
simple. Con \(t=-G_{ef}+O(j^{-2})\) y \(\hat c'_f-\hat c_f=O(j^{-2})\)
(Corolario~\ref{cor:dos-casillas}), el resto físico dividido
por \(p_ep_f\) es \(\eta+O(j^{-3})\), donde
\begin{align*}
 \eta={}&(t+G_{ef})+(\hat c'_f-\hat c_f)+\hat c_e\hat c_f+t(\hat c_e+\hat c_f)\\
        &-\mathsf E_e-\mathsf E'_f+(p_e+p_f)G_{ef}\\
     ={}&B_{\rm phys}+D^2/2+O(j^{-3}),\\
 B_{\rm phys}={}&u-D+(p_e+p_f)G_{ef}.
\end{align*}
Estos siete sumandos valen en cada estado y orientación. El cuadrado \(D^2/2\) sólo requiere
el término lineal de la inserción. Para identificar el resto \(B_{\rm phys}\) hay que comparar la
acción y el prefactor en la máscara y en su menor doble. Para hacerlo en una misma familia
retiramos sólo \(e,f\) y conservamos sus cuatro líneas. Los dos márgenes de la casilla \(i\)
pasan a ser \(1-x_i\), y todos los demás valen uno. Sean \(F_-(x),\ell_-(x)\) la acción
estacionaria y la corrección de torres de esa máscara perforada. Ponemos
\(R(x)=F_-(x)+2\sum_i(1-x_i)\log(1-x_i)\) y \(A(x)=\ell_-(x)+\sum_i\log(1-x_i)\).
Ambas expresiones regularizadas llegan a \(x=(1,1)\): el sistema visible es el menor doble,
y la masa total de su Hessiana conserva la constante de tamaño correcta.
Sean \(p=(p_e,p_f)^\top\), \(\mathbf1=(1,1)^\top\) y
\(D_p=\operatorname{diag}(p_i(1-p_i))\).
Para la Hessiana perforada de márgenes \(T_-\) y la incidencia \(a_i\),
\(R''_{il}=2\delta_{il}/(1-x_i)-a_i^\top T_-^+a_l\).
Reinsertar las casillas añade \(\sum_i[p_i(1-p_i)/(1-2p_i)]a_i a_i^\top\).
El lema del determinante en el complemento del calibre y la estacionariedad dan
\begin{align*}
 u={}&R(\mathbf1)-R(p)-(\mathbf1-p)^\top\nabla R(p)\\
     &+A(\mathbf1)-A(p)+\tfrac12\log\det(I-R''(p)D_p).
\end{align*}
Sean \(\partial=\partial_e+\partial_f\), \(\mathcal H_{\rm par}=R''(0)\) y \(\Delta C=C(\mathbf1)-C(0)\).
La derivada de la ecuación estacionaria en los dos pesos da la respuesta
\((D_p^{-1}-R''(p))^{-1}\). Se desarrollan el determinante y las respuestas de
\(\hat c_e,\hat c_f,G_{ef}\). Los términos lineales en \(p_i (\mathcal H_{\rm par})_{il}\) se cancelan por
\(\mathbf1^\top \mathcal H_{\rm par}p=\sum_i p_i(\mathcal H_{\rm par})_{ii}+(p_e+p_f)(\mathcal H_{\rm
 par})_{ef}\). Queda
\[
 B_{\rm phys}=\Delta R-\partial R(0)-\partial^2 R(0)/2+\Delta A-\partial A(0)+O_b(j^{-3}).
\]
Para obtener los coeficientes, a cada línea marcada se le añaden dos copias de signos opuestos.
En la copia positiva \(\lambda_i=1-x_i\), y en la parte ordinaria y las copias negativas
\(\lambda_i=1\). Con masa auxiliar \(s\), la medida
de número es \(M=a+s\rho\) y la de margen \(N=\lambda M\), con \(\rho=0,-1,+1\) en la parte
ordinaria, las copias negativas y las positivas. En \(s=1/j\) las copias negativas eliminan
los originales exactamente. Para la incidencia binaria \(Y_{il}\) y sus factores de
Sinkhorn~\cite{sk}
\(W_{il}\), las ecuaciones normalizadas son
\[
 w_{il}(1-s\lambda_i\lambda_l w_{il})=W_{il}\exp(z_i+y_l),\qquad
 \int Y_{il}w_{il}\,dN_C(l)=\int Y_{il}w_{il}\,dN_R(i)=1 .
\]
Denotemos por \(R_q,A_q\) los coeficientes de \(s^q\) de la acción y del prefactor regularizados
sobre \(x_e=x_f=t\). En la acción, el coeficiente de \(s^0\) es afín y \(R_1\) tiene grado
a lo sumo dos en \(t\). En el prefactor \(A_0\) es constante. Los dos términos cuadráticos
del factor local y de la diagonal de márgenes se cancelan en cada extremo marcado, y \(A_1\)
es afín. Para controlar el resto a corte fijo, sea \(\centrado=YW-1\) el núcleo basal centrado de
la máscara perforada en la buena aproximación. La brecha \(\|\centrado\|_{2\to2}\le r_b<1\)
acota \((I-\centrado^*\centrado)^{-1}\). Las cotas cuadráticas de las secciones de \(\centrado\) acotan
 también
en norma de supremo la inversa del jacobiano de márgenes en el calibre. En las copias bastan sus
ecuaciones con diagonal uno. Para \(|x_i|\le2\), tanto \(\lambda\) como la variación total
de \(m=\lambda\rho\) están acotadas. La ecuación implícita da por ello una rama holomorfa
en un disco \(|s|<\varepsilon_b\) uniforme en los pares. En el mismo disco reducido, los
grados de la Hessiana y su masa siguen cerca de uno. La serie de Fredholm del determinante
centrado admite el mayorante de Hadamard. En el origen la brecha separa su determinante de
cero, pues la suma de cuadrados singulares es acotada. Si \(M_b\) domina el determinante
en el disco inicial y \(m_b>0\) minora su módulo en cero, Cauchy acota su derivada
en el disco mitad por \(2M_b/\varepsilon_b\). En el radio reducido
\(r_b'=\min(\varepsilon_b/2,m_b\varepsilon_b/(4M_b))\) se mantiene el módulo
\(\ge m_b/2\), lo que fija la rama basal del logaritmo. Los factores \(s^{-1},s^{-2}\)
requieren uno y dos coeficientes adicionales en sus numeradores. Para una función regularizada
acotada por \(K_b\) en \(|s|\le r<r_b'\), la cola desde grado tres, evaluada en \(1/j\),
es a lo sumo \(K_b(rj)^{-3}/(1-(rj)^{-1})\). Cauchy también en el parámetro
\(t\), en un disco interior a \(|t|<2\) que contenga \([0,1]\), da la misma tasa para
sus tres primeras derivadas. Todos estos restos son \(O_b(j^{-3})\).
Con primas para las derivadas en \(t\), se obtiene
\begin{align}
 L_s&=A'_1(0)+\tfrac12R''_1(0),\notag\\
 B_{2,s}&=\tfrac12\int_0^1(1-t)^2R'''_2(t)\,dt+\int_0^1(1-t)A''_2(t)\,dt .\label{eq:estado-def}
\end{align}
En \(L_s,B_{2,s}\) el subíndice designa el estado de los cruces. Los polinomios
\(R_1,R_2,A_1,A_2\) tienen grados a lo sumo \(4,6,6,8\), de modo que nueve nodos distintos
determinan sus derivadas algebraicamente. La evaluación en punto flotante requiere su propio
control de error. Damos las recurrencias en el perfil para cada orientación y estado.
El apéndice~\ref{app:transfer} justifica su uso al promediar la máscara física.

\subsection{Datos del par y recurrencia de la acción}\label{sec:par-accion}

Llamemos \(z_{ij}=(r+ih,c+jh)\), \(i,j\in\{0,1\}\), a las cuatro casillas del cuadrado
geométrico. Numeramos las filas raíz por \((r_0,r_1)=(r,r+h)\). En la orientación \(D\)
numeramos las columnas raíz por \((c_0,c_1)=(c,c+h)\) y en la orientación \(A\) por
\((c_0,c_1)=(c+h,c)\). Así se perforan siempre \((r_0,c_0),(r_1,c_1)\). En el orden
de raíces, el bloque de incidencias es \(\left(\begin{smallmatrix}0&X\\Y&0\end{smallmatrix}\right)\).
En coordenadas del tablero esto da la siguiente tabla.
\begin{center}
\begin{tabular}{@{}ccc@{}}
\toprule
Estado \((X,Y)\) & Cruces abiertos para \(D\) & Cruces abiertos para \(A\)\\
\midrule
\((0,0)\)& ninguno & ninguno\\
\((1,0)\)& \(z_{01}\) & \(z_{00}\)\\
\((0,1)\)& \(z_{10}\) & \(z_{11}\)\\
\((1,1)\)& \(z_{01},z_{10}\) & \(z_{00},z_{11}\)\\
\bottomrule
\end{tabular}
\end{center}
Cada fila de la tabla lleva el peso \(w_{o,b}q_{o,(X,Y)}\) del apéndice~\ref{app:boundary}.
Para construir la integral a corte \(b>0\), un tipo ordinario de fila es \((u,\xi)\),
donde \(u\in[0,1]\) y \(\xi=(\xi_0,\xi_1)\) indica las aristas que llegan a las dos columnas
raíz. Escribamos \(e_1(v)=v,\ e_0(v)=1-v\). Las medidas ordinarias de fila y columna son
\[
 d\alpha_R(u,\xi)=du\prod_{k=0}^1 e_{\xi_k}(v_b(u,c_k)),\qquad
 d\alpha_C(v,\zeta)=dv\prod_{k=0}^1 e_{\zeta_k}(v_b(r_k,v)).
\]
A cada lado se añaden dos copias de cada raíz, de masa basal cero. La medida firmada
\(\rho_\sigma\), en el lado \(\sigma=R,C\), asigna masas \(-1,+1\) a sus copias
negativa y positiva. La capacidad \(\lambda\) vale \(1-t\) en las copias positivas y uno
en los demás tipos. Fijamos \(m_\sigma=\lambda\rho_\sigma\),
\(M_\sigma=\alpha_\sigma+s\rho_\sigma\) y \(N_\sigma=\alpha_\sigma+s m_\sigma\).

Sea \(V_{il}\) la incidencia condicional entre un tipo fila \(i\) y un tipo columna \(l\).
Entre tipos ordinarios de posiciones \(u,v\) vale \(v_b(u,v)\). Entre una fila ordinaria y
una copia de la columna raíz \(k\) vale \(\xi_k\). Entre una copia de la fila raíz \(k\) y
una columna ordinaria vale \(\zeta_k\). Entre copias se usa la matriz del estado indicado.
En todos los bloques \(W=W_b^0=p_b/v_b\) se evalúa en las posiciones de ambos tipos.
Con \(q_{il}=\lambda_i\lambda_l,\ \kappa_{1,il}=q_{il}W_{il}\), sustituimos la incidencia por \(V\)
en las ecuaciones de márgenes de este apéndice. Para el determinante, la corrección por
la ida y vuelta en una misma arista se calcula en la sección~\ref{sec:par-prefactor}.

\paragraph{Potenciales y densidades por orden.}
Sean \(\mathcal Z_d=(Z_d,H_d)=[s^d](z,y)\) y \(h_d(i,l)=Z_d(i)+H_d(l)\).
Partimos de \(\mathcal Z_0=0,\ w_0=W\). Para \(d=1,2,3\), el término conocido \(v_d\) de
\(w_d=[s^d]w=Wh_d+v_d\) se obtiene de \(w(1-sqw)=W\exp(z+y)\).
Omitiendo los índices de arista,
\begin{align*}
 v_1/W&=\kappa_1,\\
 v_2/W&=h_1^2/2+2\kappa_1 h_1+2\kappa_1^2,\\
 v_3/W&=h_1h_2+h_1^3/6+\kappa_1(2h_2+2h_1^2)+6\kappa_1^2h_1+5\kappa_1^3 .
\end{align*}
A orden \(d\), los dos márgenes y el calibre determinan \(\mathcal Z_d\) mediante el sistema
\begin{align*}
 Z_d(i)+\int VW H_d(l)\,d\alpha_C(l)
   &=-\int Vv_d\,d\alpha_C-\int Vw_{d-1}\,dm_C ,\\
 H_d(l)+\int VW Z_d(i)\,d\alpha_R(i)
   &=-\int Vv_d\,d\alpha_R-\int Vw_{d-1}\,dm_R ,\\
 \int Z_d\,d\alpha_R&=\int H_d\,d\alpha_C .
\end{align*}
Se resuelve primero el bloque ordinario. La brecha de \(p_b\) da la inversa en el calibre,
y las incógnitas de las copias se recuperan de su ecuación con diagonal uno.
Los lados derechos son compatibles porque las sumas de los dos márgenes respecto
de \(N_R,N_C\) coinciden y estas medidas tienen la misma masa.

\paragraph{Extracción de la acción.}
Escribimos \(\langle\alpha,\mathcal Z_d\rangle=\int Z_d\,d\alpha_R+\int H_d\,d\alpha_C\)
y análogamente para \(m\). Sean \(D_j=[s^j](N_R\otimes N_C)\). En particular,
\(D_0=\alpha_R\otimes\alpha_C,\ D_1=m_R\otimes\alpha_C+\alpha_R\otimes m_C\) y
\(D_2=m_R\otimes m_C\). En cada arista se calculan
\[
 f_1=\tfrac12 qW^2,\qquad f_2=qWw_1+\tfrac13q^2W^3,\qquad
 f_3=q(Ww_2+w_1^2/2)+q^2W^2w_1+q^3W^4/4 .
\]
En la acción dual, extraer los términos de grados uno y dos da
\begin{align*}
 R_1(t)&=-\langle\alpha,\mathcal Z_2\rangle-\langle m,\mathcal Z_1\rangle
          -\int Vf_2\,dD_0-\int Vf_1\,dD_1,\\
 R_2(t)&=-\langle\alpha,\mathcal Z_3\rangle-\langle m,\mathcal Z_2\rangle
          -\int Vf_3\,dD_0-\int Vf_2\,dD_1-\int Vf_1\,dD_2 .
\end{align*}
En efecto, \(f_1,f_2,f_3\) son los tres primeros coeficientes de
\(\sum_{a\ge2}s^{a-2}q^{a-1}w^a/a\), que se integra con signo negativo contra
\(V\,d(N_R\otimes N_C)\). La contribución de los potenciales es
\(-s^{-1}(\int z\,dN_R+\int y\,dN_C)\), módulo los términos afines en \(t\)
que se anulan en el resto de Taylor.

\subsection{Recurrencia del prefactor y traza de una arista}\label{sec:par-prefactor}

Con los \(w_d\) de la sección~\ref{sec:par-accion}, sean \(p=sqw\) y
\(\ell=w(1-p)/(1-2p)\). El grado de una fila es
\(d_R(i)=\int V_{il}\ell_{il}\,dN_C(l)\) y el de una columna se define transponiendo
la fórmula. Ambos valen uno en \(s=0\). La masa total
\(g=\int d_R\,dN_R=\int d_C\,dN_C\) también empieza en uno. Para calcular
\(d_\sigma=1+\sum_{a\ge1}s^a d_{\sigma,a}\), desarrollamos
\(\ell-w=sqw^2/(1-2sqw)\). Con \(\kappa_1=qW\), sus tres primeros coeficientes son
\begin{align*}
 E_1&=W\kappa_1,\qquad E_2=2\kappa_1 w_1+2W\kappa_1^2,\\
 E_3&=q(2Ww_2+w_1^2)+6\kappa_1^2w_1+4W\kappa_1^3 .
\end{align*}
Así \(\ell_0=W,\ \ell_1=w_1+E_1,\ \ell_2=w_2+E_2\). Con \(E_0=0\), los márgenes
\(\int Vw\,dN_{\rm opuesto}=1\) dan para \(a=1,2,3\)
\[
 d_{R,a}=\int VE_a\,d\alpha_C+\int VE_{a-1}\,dm_C,\qquad
 d_{C,a}=\int VE_a\,d\alpha_R+\int VE_{a-1}\,dm_R .
\]
Ponemos \(g_a=\int d_{R,a}\,d\alpha_R+\int d_{R,a-1}\,dm_R\) para \(a=1,2\),
con \(d_{R,0}=1\). Para cada lado, los coeficientes \(D_{\sigma,a}^{\log}=[s^a]\log d_\sigma\)
son \(d_1,\ d_2-d_1^2/2,\ d_3-d_1d_2+d_1^3/3\) en órdenes uno, dos y tres,
respectivamente, con \(d_a=d_{\sigma,a}\).

\paragraph{Términos locales y determinante.}
Escribamos \(\mu_a=[s^a](M_R\otimes M_C)\), con \(\mu_a=0\) para \(a\notin\{0,1,2\}\).
Para \(\kappa_1=qW,\ \kappa_2=qw_1,\ \kappa_3=qw_2\), los coeficientes de
\(-\log(1-2p)+2\log(1-p)\) en órdenes dos, tres y cuatro son
\[
 C_2=\kappa_1^2,\qquad C_3=2\kappa_1\kappa_2+2\kappa_1^3,\qquad
 C_4=\kappa_2^2+2\kappa_1\kappa_3+6\kappa_1^2\kappa_2+\tfrac72\kappa_1^4 .
\]
Para \(a=1,2\), las aristas, la masa total y las diagonales de la Hessiana aportan
\begin{align*}
 H_a^{\rm loc}={}&\tfrac12\sum_{j=2}^{a+2}\int V C_j\,d\mu_{a+2-j}
 +\tfrac12[s^a]\log(1+sg_1+s^2g_2)\\
 &-\tfrac12\sum_{\sigma=R,C}
 \left(\int D_{\sigma,a+1}^{\log}\,d\alpha_\sigma+
       \int D_{\sigma,a}^{\log}\,d\rho_\sigma\right).
\end{align*}
Las distintas medidas proceden de la masa de la Hessiana \(g\), integrada contra \(N_R\),
y de la suma de sus logaritmos de grado, integrada contra la medida de número \(M_\sigma\).

Para los bloques centrados, los operadores \(K_a,T_a\), \(a=0,1,2\), se obtienen
extrayendo el coeficiente de cada núcleo junto con su medida de integración.
\begin{align*}
 K_a(i,dl)&=[s^a]\left\{\left(\frac{V_{il}\ell_{il}}{d_R(i)}-\frac{d_C(l)}g\right)
                       (d\alpha_C(l)+s\,dm_C(l))\right\},\\
 T_a(l,di)&=[s^a]\left\{\left(\frac{V_{il}\ell_{il}}{d_C(l)}-\frac{d_R(i)}g\right)
                       (d\alpha_R(i)+s\,dm_R(i))\right\}.
\end{align*}
Para evaluar los inversos escalares
se usa \((1+sa+s^2b)^{-1}=1-sa+s^2(a^2-b)+O(s^3)\). Sean
\[
 I_0=(I-K_0T_0)^{-1},\qquad X_1=K_1T_0+K_0T_1,\qquad
 X_2=K_2T_0+K_1T_1+K_0T_2 .
\]
El desarrollo de \(-\tfrac12\log\det(I-KT)\) para los dos bloques de medias da
\[
 H_1^{\rm tr}=\tfrac12\operatorname{tr}(I_0X_1),\qquad
 H_2^{\rm tr}=\tfrac12\operatorname{tr}(I_0X_2)+\tfrac14\operatorname{tr}((I_0X_1)^2).
\]

\paragraph{La ida y la vuelta por una casilla ordinaria.}
En una traza de longitud dos figura el mismo indicador \(I\) en ambas direcciones.
Para una arista entre tipos ordinarios, \(\mathbb E[I^2]=V\) y
\(\mathbb E[I]^2=V^2\), condicionalmente a los tipos. Al multiplicar los dos núcleos
centrados, los términos constantes y lineales tienen las medias correctas. En el
término cuadrático falta \(V(1-V)\ell^2/(d_Rd_C)\). Sea \(d\mu_*=\tfrac12
V(1-V)\,d\alpha_R\,d\alpha_C\), incluido el peso \(1/2\) de la traza. Para
\(r=d_{R,1}, c=d_{C,1}, r_2=d_{R,2}, c_2=d_{C,2}\), las dos contribuciones son
\begin{align*}
 H_1^{\rm rep}&=\int \big[2W\ell_1-(r+c)W^2\big]\,d\mu_*,\\
 H_2^{\rm rep}&=\int \big[\ell_1^2+2W\ell_2-2(r+c)W\ell_1
              +(r^2+c^2+rc-r_2-c_2)W^2\big]\,d\mu_* .
\end{align*}
Los tipos raíz tienen incidencia fijada. Todas sus repeticiones usan ese mismo bit.
Concluimos que, para cada orientación y estado, el prefactor pedido es
\[
 A_a(t)=H_a^{\rm loc}+H_a^{\rm tr}+H_a^{\rm rep},\qquad a=1,2 .
\]

\paragraph{Ejemplo de un sumando de \(A_2(1)\).}
Tomemos como dato de prueba el núcleo constante \(v=1/2\), cuyo escalado es \(p_b=1\).
Fijemos el estado \(X=Y=0\) y la capacidad \(1-t=0\). En ambos lados los patrones
ordinarios \(00,10,01,11\) tienen masa \(1/4\), y \(W=2\). Consideremos la ida y vuelta
entre una fila de patrón \(00\) y una columna también de patrón \(00\).
Para comprobar las entradas de la recurrencia, por simetría \(Z_a=H_a\). En el orden de
los cuatro patrones seguido de las dos copias negativas y las dos positivas, los sistemas dan
\[
 Z_1=(-2,0,0,2,-3,-3,-1,-1),\qquad
 Z_2=(2,-6,-6,-10,11/2,11/2,15/2,15/2).
\]
En la arista elegida resultan \(w_1=-4,\ w_2=8,\ \ell_1=0,\ \ell_2=8\).
En ambos extremos \(d_1=2,\ d_2=8\). De aquí
\[
 [s^2]\frac{\ell^2}{d_Rd_C}
 =[s^2]\frac{(2+8s^2)^2}{(1+2s+8s^2)^2}=16 .
\]
El sumando de \(H_2^{\rm rep}\) para estos dos patrones es, pues,
\(\tfrac12(1/4)^2(1/2)(1-1/2)\,16=1/8\).
La integración en las dos posiciones tiene masa uno porque el dato de prueba es constante.


\section{De las máscaras marcadas al perfil}\label{app:transfer}
Dos casillas pueden depender de la supervivencia de una misma línea. Tratar sus indicadores
como independientes cambiaría la media del par, incluso si el perfil de cada casilla fuese
exacto. Por ello transferimos la contribución ponderada de cada estado con su propia masa.
Para un patrón fijo identificamos las etiquetas afines iguales y tomamos cada indicador
distinto una sola vez. Con \(d_*\) índices libres,
una coincidencia accidental satisface una ecuación afín no constante adicional y cuesta
\(O(n^{-1})\) al dividir por \(n^{d_*}\). En la trayectoria regular auxiliar, a \(x\ge b\), los
riesgos de las líneas axiales y diagonales vivas son \(1/m\) y \(\ell_f/m\), respectivamente.
Los segundos distan de sus valores deterministas a lo sumo \(\eta/n\) a precisión fija.
Dos líneas distintas se ocupan simultáneamente con probabilidad \(O_b(n^{-2})\). Bonferroni da
el producto de sus supervivencias de un paso con error \(O_{b,v}(\eta/n+n^{-2})\).
Multiplicamos cada indicador actual por su supervivencia determinista futura hasta el corte.
Una pieza altera \(O_v(n^{d_*-1})\) sumandos. Azuma para la media centrada, con saltos
\(O_{b,v}(1/n)\), y el control de regularidad dan una cola exponencial a precisión fija.
La unión sobre líneas marcadas es polinómica. Los cocientes de densidades hasta la parada
preservan la cola, incluida la cota \(\exp(O_b(n^{2/3}))\) para \(X\) al cambiar a \(\mu\).
Así, en la media normalizada cada etiqueta axial viva aporta \(b\) y cada etiqueta diagonal
distinta aporta \(U_f(b)\).

En particular, las sumas por línea viva de \(YW_b^0/j\) tienen márgenes \(1+o_{\Pr}(1)\)
uniformemente y dos columnas vivas tienen \(\ge c_bj\) vecinos comunes. En potenciales
logarítmicos de columna, el mapa que normaliza primero filas y después columnas tiene
derivada estocástica de entradas \(\ge c'_b/j\) en la bola de oscilación unidad. Su
contracción en oscilación da \(jD_{il}=Y_{il}W_b^0(i,l)(1+o_{\Pr}(1))\) para el escalado
de Sinkhorn, uniformemente en el soporte. Al retirar las dos casillas distinguidas, cada
margen pierde \(O_b(1/j)\) y dos columnas pierden a lo sumo dos vecinos comunes. La misma
prueba sigue aplicándose. La ocupación estacionaria de torres cumple
\(p_{il}/D_{il}=1+O_b(1/j)\) en el soporte. En los patrones integrados, las incidencias desde una
fila adicional \(y\) hacia las columnas marcadas tienen ley producto con parámetros
\(v_b(y,c_0),v_b(y,c_1)\). Las distintas columnas dan etiquetas diagonales distintas en cada
familia, salvo los \(O(1)\) valores de \(y\) que coinciden con una etiqueta distinguida.
Se puede multiplicar la media por cualquiera de los cuatro indicadores acotados de estado.
Para el término cuadrático de \(\eta\), sean \(B=p_b-1\),
\(G_R=BB^*, T_R=G_R(I-G_R)^{-1}, U=T_RB\), con los operadores transpuestos en columnas.
Sean \(d_R=\operatorname{diag}(T_R-G_R)\) y su análogo \(d_C\), con
\(\bar d=\int d_R=\int d_C\). Definimos
\(\psi_R=2[(B+U)d_C-T_Rd_R]-d_R\) y su versión de columnas. En los cuatro extremos
marcados (dos filas \(r_0,r_1\) y dos columnas \(c_0,c_1\)), las identidades de derivadas dan
\begin{align*}
 R''_1(0)/2={}&XW^0_{01}+YW^0_{10}+\sum_{\square} U-T_R(r_0,r_1)-T_C(c_0,c_1)
              -\tfrac12\sum_* d,\\
 A'_1(0)={}&-1+2\bar d+\tfrac12\sum_* \psi .
\end{align*}
El asterisco recorre los cuatro extremos y el cuadrado sus cuatro casillas; se ordenan las columnas
para que los casillas distinguidas sean \(00,11\). Para la primera identidad se derivan los
márgenes en el parámetro de ocupación. En los tipos ordinarios la respuesta media resuelve
\((I+N)y=N\alpha\), donde \(N\) tiene bloques cruzados \(p_b,p_b^*\) y \(\alpha\) inserta
los dos extremos de cada lado. Los bloques centrados de \(N(I+N)^{-1}N\) son \(T_R,-U,-U^*,T_C\).
En cada vértice marcado una incidencia repetida da su varianza de Bernoulli. Para el prefactor
se derivan el factor local, las diagonales de márgenes, la masa total y el determinante centrado.
Su suma es \(1+\langle d-1,y\rangle-\sum_*d/2\); sustituir la solución centrada y su modo
constante da la segunda identidad. La parte lineal límite \(jD\) es la suma \(L_s\)
en cada submedida. Al elevar al cuadrado resulta \(L_s^2/2\).

El argumento de patrones sirve a cualquier longitud fija, pero el segundo coeficiente contiene
la inversa de márgenes y, por tanto, recorridos de longitud arbitraria. Truncamos su expansión
en recorridos alternantes centrados. Cada recorrido de longitud \(q\ge2\) entre extremos fijos
está acotado por
\(C_b^2r_b^{q-2}\), por las dos normas de sección y la brecha \(r_b<1\). A orden dos en el
logaritmo del determinante hay a lo sumo dos inserciones en cada traza; su cola está acotada
por \(C_b\sum_{q>N}(1+q)^2r_b^{\max(2q-C,0)}\). A truncación fija, las recurrencias de márgenes
adjuntan árboles enraizados a una casilla o a una traza alternante. Las coincidencias de
etiquetas libres distintas cuestan \(O_{b,N}(1/n)\). Se conserva cada indicador repetido antes de
tomar la media. Una traza de longitud dos en una arista libre de media condicional \(V\) usa
\(Y^2=Y\) a la ida y vuelta. El exceso respecto al producto de medias aporta
\(V(1-V)\ell^2/(2d_i d_l)\), con \(\ell\) el peso normalizado de la arista en la Hessiana y
\(d_i,d_l\) sus grados de fila y columna. Las aristas raíz ya tienen fijados sus indicadores.
Así la varianza ordinaria se cuenta una vez. Primero \(n\to\infty\) a truncaciones fijas;
después se retiran las truncaciones. Sólo entonces \(b\downarrow0\): la positividad de \(v_b\)
da una minoración uniforme de \(p_b^*p_b\) cerca de cero y convergencia dominada en cada integral.
Con \(\mathcal I_0(C)=\sum_o\int_\Omega w_{o,0}\sum_s q_{o,s,0}C_{o,s}\), se obtiene
\[
 \lim_{b\downarrow0}\limsup_n\left|
 n\mathbb E_\mu[\textstyle\sum_{\{e,f\}}p_ep_f B_{\rm phys}(e,f);F_b]
 -\mathcal I_0(B_{2,s})\right|=0. \tag{F2}
\]
Así la frontera firmada por pares es \(J=-B_2+L_2\), con \(B_2=\mathcal I_0(B_{2,s})\)
y \(L_2=-\mathcal I_0(L_s^2/2)\). En el estado \(11\) cada orientación aporta su coeficiente
de una línea y la interacción de las dos orientaciones aparece una sola vez en \(-S\).
Sumar \(J\) a las otras dos medias da la constante \(P_T+J-S\) que entra en la comparación
de cortes, una vez tomado el límite en \(n\) para cada corte fijo.

\section{Comparación de dos cortes e identificación del límite}\label{app:closing}
Las estimaciones anteriores valen para cada corte admisible fijo. Sus umbrales de \(n\)
pueden depender del corte, por lo que aún hay que justificar el paso al extremo.
La cantidad que aproximan en el tablero inicial es la misma para todos los cortes.
Sean \(a_n=nR_Q(n)\), \(a_*=F(1)-1/12\), \(H=P_T+J-S\), y \(b>0\) un corte
admisible. Existen cortes admisibles arbitrariamente próximos a cero. Definimos a corte positivo
\(I(b)=\int_b^1 k(x)x^{-2}\,dx\) y \(\Phi(b)=a_*+I(b)+H\).
Las tres identificaciones de frontera dan \(\eta(b)=\limsup_n|B_n(b)-H|\to0\).
La identidad interior vale en \([b,1]\). Al combinarla con \eqref{eq:signed-mean},
la expansión de frontera con error \(O(n/j^2)\) y \(n g_T(M_n)\to a_*\), queda
\[
 \limsup_n|a_n-\Phi(b)|\le C_\delta b+\eta(b)=:\varepsilon(b).
\]
\begin{lema}[Comparación de dos cortes]\label{lem:dos-cortes}
Sea \((a_n)\) una sucesión real y, para cada corte admisible \(b\), sea \(\Phi(b)\in\mathbb R\) con
\(\limsup_n|a_n-\Phi(b)|\le\varepsilon(b)\), donde \(\varepsilon(b)\to0\) al bajar el corte. Entonces
\(\Phi(b)\) tiene límite finito \(\Phi_0\) a lo largo de los cortes admisibles y \(a_n\to\Phi_0\).
\end{lema}
\begin{proof}
Para dos cortes admisibles fijos comparamos \(\Phi(b),\Phi(b')\) con el mismo \(a_n\), tomando \(n\)
suficientemente grande para ambas estimaciones. La desigualdad triangular y el límite superior
dan \(|\Phi(b)-\Phi(b')|\le\varepsilon(b)+\varepsilon(b')\). Se cumple el criterio de Cauchy a lo
largo de los cortes admisibles y \(\Phi(b)\to\Phi_0\in\mathbb R\). Para cada corte,
\(\limsup_n|a_n-\Phi_0|\le\varepsilon(b)+|\Phi(b)-\Phi_0|\). Como el lado izquierdo es independiente
de \(b\), al bajar el corte resulta \(a_n\to\Phi_0\).
\end{proof}
En particular \(I(b)\) tiene límite finito a lo largo de dichos cortes. En todo este argumento
el perfil físico, las leyes detenidas y sus estimaciones se han usado a corte positivo fijo.

\begin{proof}[Final de la demostración del Teorema~\ref{teo:principal} y del
 Corolario~\ref{cor:derivada}]
Para identificar el límite con la integral impropia ordinaria queda decidir si \(k(x)/x^2\)
puede tener un término proporcional a \(1/x\). La identidad de orden cero del perfil da
\(k(0)=0\). La resolvente proyectada de líneas es analítica
cerca de cero: el bloque de Schur diagonal tiene la forma \(\Lambda_u+x\mathcal N\), la vacancia
tiene mínimo positivo y el bloque axial tiene brecha espectral. La inversa y las secciones de su
núcleo permanecen acotadas en discos complejos interiores. La integración dominada en las
contracciones finitas de las densidades da así la analiticidad de \(k,G\). Esta regularidad se
obtiene para cualquier valor que pueda tener \(k'(0)\). En un intervalo suficientemente
pequeño \([0,x_0]\), escribimos
\(k(x)=k'(0)x+x^2 h(x)\), con \(h\) analítica y acotada. Para \(0<b<x_0\),
\[
 I(b)=I(x_0)+k'(0)\log(x_0/b)+\int_b^{x_0}h(x)\,dx .
\]
La última integral permanece acotada. La convergencia de \(I(b)\) por los cortes admisibles
fuerza \(k'(0)=0\). De \(k'=d'-xG''\) se deduce \(d'(0)=0\) para el perfil físico.
La singularidad del integrando es entonces evitable y la integral impropia converge
absolutamente. Por consiguiente, el valor de \(\Phi_0\) es
\[
 \lim_{n\to\infty}n\{\log Q(n)-\log Q_{\rm pred}(n)\}
   =F(1)-\tfrac1{12}+\int_0^1\frac{k(x)}{x^2}\,dx+P_T+J-S .
\]
\end{proof}

